% =============================================================================
% main.tex — Esqueleto da Monografia (abntex2)
% Chronicles Of The Lost Word: Integração entre Proveniência e
% Aprendizado por Reforço para Comportamento de Chefe em Jogo 2D
% =============================================================================
\documentclass[
    12pt,                % Tamanho de fonte
    openright,           % Capítulos começam em página ímpar
    oneside,             % Para impressão em um lado só (twoside para frente e verso)
    a4paper,             % Tamanho do papel
    chapter=TITLE,       % Títulos de capítulo em maiúsculas
    section=TITLE,       % Títulos de seção em maiúsculas
    english,             % Idioma secundário (para abstract em inglês)
    brazil               % Idioma principal
]{abntex2}

% ---------------------------------------------------------------------------
% PACOTES ESSENCIAIS
% ---------------------------------------------------------------------------

% --- Codificação e fontes ---
\usepackage[T1]{fontenc}        % Codificação de saída
\usepackage[utf8]{inputenc}     % Codificação de entrada (UTF-8)
\usepackage{lmodern}            % Fonte Latin Modern (boa para PDF)

% --- Português / Babel ---
\usepackage[brazil]{babel}      % Hifenização e strings em português

% --- Formatação ABNT ---
\usepackage{indentfirst}        % Indenta o primeiro parágrafo de cada seção
\usepackage{microtype}          % Melhora a tipografia geral

% --- Matemática ---
\usepackage{amsmath}            % Ambientes matemáticos
\usepackage{amssymb}            % Símbolos matemáticos
\usepackage{amsfonts}           % Fontes matemáticas

% --- Figuras e gráficos ---
\usepackage{graphicx}           % Inclusão de imagens
\usepackage{float}              % Posicionamento forçado de floats [H]
\usepackage{subcaption}         % Sub-figuras (a), (b)...
\graphicspath{{figuras/}}       % Pasta padrão de figuras

% --- Tabelas ---
\usepackage{booktabs}           % Linhas horizontais elegantes
\usepackage{multirow}           % Células multi-linha
\usepackage{array}              % Melhorias em colunas de tabela
\usepackage{longtable}          % Tabelas que quebram página

% --- Algoritmos e pseudocódigo ---
\usepackage{algorithm}                  % Ambiente flutuante para algoritmos
\usepackage[noend]{algpseudocode}       % Pseudocódigo sem "end"
\floatname{algorithm}{Algoritmo}        % Renomeia "Algorithm" para "Algoritmo"
\renewcommand{\algorithmicrequire}{\textbf{Entrada:}}
\renewcommand{\algorithmicensure}{\textbf{Saída:}}

% --- Código-fonte ---
\usepackage{listings}           % Listagens de código
\usepackage{xcolor}             % Cores personalizadas

% Configuração padrão de listings
\lstset{
    basicstyle=\ttfamily\small,
    breaklines=true,
    frame=single,
    numbers=left,
    numberstyle=\tiny\color{gray},
    keywordstyle=\color{blue},
    commentstyle=\color{green!60!black},
    stringstyle=\color{red!70!black},
    tabsize=4,
    showstringspaces=false,
    captionpos=b
}

% --- Citações e referências ---
\usepackage[brazilian,hyperpageref]{backref}
\usepackage[alf]{abntex2cite}   % Citação no padrão ABNT (autor-data)

% --- Hiperlinks ---
\usepackage{hyperref}
\hypersetup{
    pdftitle={Chronicles Of The Lost Word: Integração entre Proveniência e Aprendizado por Reforço para Comportamento de Chefe em Jogo 2D},
    pdfauthor={Matheus Oliveira},
    pdfsubject={Trabalho de Conclusão de Curso em Ciência da Computação},
    colorlinks=true,
    linkcolor=blue,
    citecolor=blue,
    filecolor=magenta,
    urlcolor=blue,
    bookmarksdepth=4
}

% --- Backref em português ---
\renewcommand{\backrefpagesname}{Citado na(s) página(s):~}
\renewcommand{\backref}{}
\renewcommand*{\backrefalt}[4]{
    \ifcase #1
        Nenhuma citação no texto.
    \or
        Citado na página #2.
    \else
        Citado #1 vezes nas páginas #2.
    \fi
}

% ---------------------------------------------------------------------------
% DADOS DO TRABALHO (CAPA E FOLHA DE ROSTO)
% ---------------------------------------------------------------------------

\titulo{Chronicles Of The Lost Word: Integração entre Proveniência e Aprendizado por Reforço para Comportamento de Chefe em Jogo 2D}

\autor{Matheus Oliveira de Lima}

\local{Nova Iguaçu -- Rio de Janeiro}

\data{2026}

\orientador{Bruno Jose Dembogurski }

% \coorientador{Prof. Dr. Nome do Coorientador}

\instituicao{%
    Universidade Federal Rural do Rio de Janeiro (UFRRJ)
    \par
    Campus Nova Iguaçu
    \par
    Curso de Ciência da Computação
}

\tipotrabalho{Trabalho de Conclusão de Curso}

\preambulo{%
    Monografia apresentada ao Curso de Ciência da Computação
    da Universidade Federal Rural do Rio de Janeiro como requisito parcial para a obtenção
    do grau de Bacharel em Ciência da Computação.
}

% ---------------------------------------------------------------------------
% INÍCIO DO DOCUMENTO
% ---------------------------------------------------------------------------
\begin{document}

% Seleciona idioma principal
\selectlanguage{brazil}

% Retira espaço extra obsoleto entre as frases
\frenchspacing

% ===========================
% ELEMENTOS PRÉ-TEXTUAIS
% ===========================
\pretextual

% --- Capa ---
\imprimircapa

% --- Folha de rosto ---
\imprimirfolhaderosto*  % O asterisco (*) gera a ficha catalográfica no verso

% --- Folha de aprovação (modelo) -> colocar depois ---
% \imprimirfolhadeaprovacao ou incluir PDF.

% --- Dedicatória ---
% \begin{dedicatoria}
%     \vspace*{\fill}
%     \centering
%     \noindent
%     \textit{Texto da dedicatória.}
%     \vspace*{\fill}
% \end{dedicatoria}

% --- Agradecimentos ---
% \begin{agradecimentos}
%     Agradecimentos aqui.
% \end{agradecimentos}

% --- Epígrafe ---
% \begin{epigrafe}
%     \vspace*{\fill}
%     \begin{citacao}
%         ``Texto da epígrafe.'' \\
%         (Autor da epígrafe)
%     \end{citacao}
% \end{epigrafe}

% --- Resumo em português ---
\begin{resumo}
    % TODO: Escrever o resumo do TCC (150 a 500 palavras).
    O desenvolvimento de comportamentos adaptativos para Inteligência Artificial (IA) em jogos digitais é um desafio recorrente, frequentemente limitado à previsibilidade das Máquinas de Estados Finitos (FSM) ou à complexidade de modelar funções de recompensa em Aprendizado por Reforço (RL). Este trabalho investiga a integração de Grafos de Proveniência como mecanismo de \textit{Reward Shaping} para o treinamento de um agente RL utilizando o algoritmo \textit{Proximal Policy Optimization} (PPO). A metodologia consistiu no desenvolvimento de um \textit{vertical slice} de um jogo 2D (\textit{Chronicles Of The Lost Word}), onde sessões de jogo humano contra um chefe FSM (\textit{baseline}) tiveram suas cadeias causais registradas estruturalmente. A partir desses grafos, extraíram-se sequências de ações vitoriosas que serviram como sinal auxiliar de recompensa para o treinamento do agente na plataforma Unity ML-Agents. A avaliação quantitativa comparou 10 sessões jogáveis contra o chefe FSM e 10 contra o modelo PPO treinado. Os resultados evidenciaram que o agente PPO desenvolveu um comportamento mensuravelmente distinto, caracterizado por maior persistência e agressividade, resultando em combates mais longos e infligindo a única derrota ao jogador na amostra, embora tenha apresentado repetição mecânica em suas ações. Conclui-se que a utilização de sequências extraídas de proveniência para guiar o aprendizado por reforço é uma abordagem viável para gerar comportamentos alternativos em jogos digitais sem a necessidade de imitação direta.

    \textbf{Palavras-chave}: Aprendizado por Reforço. Proveniência de Dados. PPO.
    Jogos Digitais. Unity ML-Agents. Reward Shaping. FSM.
\end{resumo}

% --- Resumo em inglês (Abstract) ---
\begin{resumo}[Abstract]
    \begin{otherlanguage*}{english}
        The development of adaptive behaviors for Artificial Intelligence (AI) in digital games is a recurring challenge, often limited by the predictability of Finite State Machines (FSM) or the complexity of modeling reward functions in Reinforcement Learning (RL). This study investigates the integration of Provenance Graphs as a Reward Shaping mechanism for training an RL agent using the Proximal Policy Optimization (PPO) algorithm. The methodology consisted of developing a vertical slice of a 2D game (\textit{Chronicles Of The Lost Word}), where human gameplay sessions against an FSM boss (baseline) had their causal chains structurally recorded. From these graphs, sequences of winning actions were extracted to serve as an auxiliary reward signal for training the agent on the Unity ML-Agents platform. The quantitative evaluation compared 10 playable sessions against the FSM boss and 10 against the trained PPO model. The results showed that the PPO agent developed a measurably distinct behavior, characterized by greater persistence and aggressiveness, resulting in longer combats and inflicting the only player defeat in the sample, although it presented mechanical repetition in its actions. It is concluded that using sequences extracted from provenance to guide reinforcement learning is a viable approach to generating alternative behaviors in digital games without the need for direct imitation.

        \textbf{Keywords}: Reinforcement Learning. Data Provenance. PPO.
        Digital Games. Unity ML-Agents. Reward Shaping. FSM.
    \end{otherlanguage*}
\end{resumo}

% --- Lista de figuras ---
\pdfbookmark[0]{\listfigurename}{lof}
\listoffigures*
\cleardoublepage

% --- Lista de tabelas ---
\pdfbookmark[0]{\listtablename}{lot}
\listoftables*
\cleardoublepage

% --- Lista de algoritmos (opcional) ---
% \pdfbookmark[0]{Lista de algoritmos}{loa}
% \listofalgorithms
% \cleardoublepage

% --- Lista de abreviaturas e siglas ---
\begin{siglas}
    \item[FSM]   Máquina de Estados Finitos (\textit{Finite State Machine})
    \item[IA]    Inteligência Artificial
    \item[ML]    Aprendizado de Máquina (\textit{Machine Learning})
    \item[ONNX]  Open Neural Network Exchange
    \item[PPO]   Proximal Policy Optimization
    \item[RL]    Aprendizado por Reforço (\textit{Reinforcement Learning})
    \item[TCC]   Trabalho de Conclusão de Curso
    \item[GAIL]  Generative Adversarial Imitation Learning
    \item[NPC]   Personagem Não-Jogável (\textit{Non-Player Character})
    \item[CSV]   \textit{Comma-Separated Values}
\end{siglas}

% --- Sumário ---
\pdfbookmark[0]{\contentsname}{toc}
\tableofcontents*
\cleardoublepage

% ===========================
% ELEMENTOS TEXTUAIS
% ===========================
\textual

% --- Capítulos ---
% =============================================================================
% Capítulo 1 — Introdução
% =============================================================================
\chapter{Introdução}
\label{cap:introducao}

\section{Contextualização}
\label{sec:contextualizacao}

A indústria de jogos digitais consolidou-se como uma das mais relevantes no
cenário do entretenimento contemporâneo, movimentando bilhões de dólares
anualmente e alcançando audiências que superam outras mídias tradicionais
\cite{newzoo2024}. Parte significativa dessa evolução deve-se à sofisticação
crescente dos sistemas de Inteligência Artificial (IA) empregados no controle
de Personagens Não-Jogáveis (NPCs), cuja qualidade comportamental exerce
influência direta sobre a experiência do jogador \cite{millington2019ai}.

Historicamente, a IA em jogos comerciais apoiou-se em técnicas determinísticas,
como Máquinas de Estados Finitos (FSM) e Árvores de Comportamento
(\textit{Behavior Trees}). Essas abordagens oferecem previsibilidade e
facilidade de depuração, qualidades fundamentais em contextos de produção.
Contudo, comportamentos baseados em regras fixas tendem a se tornar
previsíveis após poucas interações, reduzindo o desafio percebido e
comprometendo a rejogabilidade \cite{yannakakis2018artificial}.

O Aprendizado por Reforço (\textit{Reinforcement Learning} --- RL) surgiu como
alternativa promissora para a criação de agentes adaptativos, capazes de
ajustar suas estratégias por meio de interações com o ambiente
\cite{sutton2018reinforcement}. Plataformas como o Unity ML-Agents
\cite{juliani2020unity} democratizaram o acesso a algoritmos de RL no contexto
de motores de jogo, permitindo que desenvolvedores treinem NPCs diretamente
em ambientes simulados. Em particular, o algoritmo \textit{Proximal Policy
Optimization} (PPO) \cite{schulman2017ppo} destacou-se por sua estabilidade
durante o treinamento e por seu bom desempenho em espaços de ação discretos,
características que o tornam adequado para o controle de personagens em jogos
2D.

Entretanto, a adoção de RL em jogos não está isenta de desafios. Agentes
treinados exclusivamente por tentativa e erro podem desenvolver
comportamentos inesperados, difíceis de interpretar e de ajustar por
projetistas de jogo. Nesse contexto, a busca por mecanismos que tornem o
processo de aprendizado mais interpretável --- isto é, que permitam
compreender \textit{por que} o agente aprendeu determinado comportamento ---
constitui uma questão aberta e relevante para a área.

\section{Problema de Pesquisa}
\label{sec:problema}

A definição de funções de recompensa constitui um dos aspectos mais críticos
do treinamento por RL. Quando projetadas manualmente, essas funções tendem a
ser exaustivas, frágeis e difíceis de generalizar: pequenas alterações no
\textit{design} do jogo podem exigir reajustes significativos no esquema de
recompensas \cite{ng1999reward}. Em cenários complexos, a engenharia manual
de recompensas torna-se um gargalo de desenvolvimento.

Uma alternativa amplamente explorada na literatura é a aprendizagem por
imitação. Técnicas como \textit{Generative Adversarial Imitation Learning}
(GAIL) permitem que um agente aprenda a replicar o comportamento de um
jogador humano a partir de demonstrações gravadas. Embora eficaz em diversos
cenários, essa abordagem apresenta uma limitação conceitual relevante: o
agente tende a \textit{copiar} o comportamento observado, sem necessariamente
compreender a relação causal entre as ações executadas e o sucesso obtido. O
resultado é um modelo que reproduz padrões comportamentais sem internalizar a
estratégia subjacente \cite{ho2016gail}.

Emerge, portanto, uma questão de pesquisa: \textbf{é possível guiar o
aprendizado por reforço de um agente de jogo utilizando informações causais
extraídas de sessões de jogo anteriores, sem recorrer à imitação direta de
demonstrações?}

A Proveniência de Dados (\textit{Data Provenance}) oferece um caminho para
endereçar essa questão. Originalmente desenvolvida para rastreamento de
linhagem em bancos de dados e fluxos científicos, a proveniência permite
registrar cadeias causais de eventos --- não apenas \textit{o que} aconteceu,
mas \textit{por que} aconteceu, preservando relações de causa e efeito entre
ações e consequências \cite{provenance_base}. Aplicada ao contexto de jogos
digitais, a proveniência possibilita a construção de grafos que representam a
sequência causal de eventos em uma sessão de jogo, desde ações do jogador até
desfechos como vitória ou derrota \cite{kohwalter2020provenance}.

\section{A Proposta}
\label{sec:proposta}

O presente trabalho propõe uma abordagem que integra Proveniência de Dados e
Aprendizado por Reforço para o treinamento de um agente de jogo. A ideia
central é utilizar grafos de proveniência, gerados a partir de sessões jogadas
contra um chefe controlado por FSM, como fonte de informação causal para guiar
o aprendizado de um agente PPO.

O pipeline proposto opera da seguinte forma:

\begin{enumerate}
    \item Um jogador humano enfrenta o chefe baseline (FSM) em múltiplas
          sessões de jogo.
    \item O sistema de proveniência registra cada evento relevante da sessão
          --- ataques, danos, pulos, esquivas, mortes --- como nós de um
          grafo causal, preservando as relações de causa e efeito entre eles.
    \item Um processo de filtragem seleciona apenas as sessões encerradas
          com vitória do jogador e extrai, a partir dos grafos correspondentes,
          sequências de ações que precederam eventos de dano no chefe.
    \item Essas sequências vitoriosas são convertidas em um artefato
          estruturado e carregadas por um módulo de \textit{Reward Shaping}
          \cite{ng1999reward}, que concede recompensa adicional ao agente PPO
          sempre que este executa padrões de ação compatíveis com os padrões
          extraídos.
    \item O agente é treinado com PPO padrão, utilizando o
          \textit{Reward Shaping} por proveniência como sinal auxiliar. O
          algoritmo continua explorando o ambiente livremente, mas recebe
          orientação indireta das ações humanas bem-sucedidas.
\end{enumerate}

Essa abordagem distingue-se da imitação direta em um aspecto fundamental: o
agente não reproduz trajetórias gravadas nem utiliza um discriminador
adversário para aprender políticas. Em vez disso, ele recebe um sinal de
recompensa adicional que reflete padrões causais associados ao sucesso,
preservando a capacidade exploratória do PPO e mantendo a rastreabilidade das
decisões de \textit{design} do treinamento.

A contribuição acadêmica central, portanto, reside na demonstração de que
sequências extraídas de grafos de proveniência podem ser utilizadas como sinal
auxiliar de recompensa em um agente PPO, produzindo um chefe com comportamento
mensuravelmente distinto do \textit{baseline} FSM em um ambiente de jogo 2D.

\section{Objetivos}
\label{sec:objetivos}

\subsection{Objetivo Geral}
\label{subsec:objetivo-geral}

Desenvolver e avaliar um pipeline de integração entre Proveniência de Dados e
Aprendizado por Reforço (PPO) para o treinamento de um agente de chefe em um
\textit{vertical slice} de jogo 2D, utilizando sequências de ações vitoriosas
extraídas de grafos de proveniência como \textit{Reward Shaping}, e comparando
o comportamento resultante com um chefe \textit{baseline} controlado por
Máquina de Estados Finitos.

\subsection{Objetivos Específicos}
\label{subsec:objetivos-especificos}

Para alcançar o objetivo geral, foram definidos os seguintes objetivos
específicos:

\begin{enumerate}
    \item Implementar o \textit{vertical slice} do jogo \textit{Chronicles Of
          The Lost Word} em Unity, contendo uma arena de chefe com mecânicas
          jogáveis de combate 2D.

    \item Desenvolver um chefe \textit{baseline} controlado por Máquina de
          Estados Finitos (FSM), com comportamento determinístico e
          previsível, adequado para servir como referência de comparação.

    \item Projetar e implementar um sistema de proveniência capaz de registrar
          eventos de jogo como nós de um grafo causal, preservando relações
          de causa e efeito entre ações do jogador e consequências no
          ambiente.

    \item Construir um processo automatizado de extração de sequências de
          ações vitoriosas a partir dos grafos de proveniência gerados em
          sessões contra o chefe FSM.

    \item Implementar um módulo de \textit{Reward Shaping} que carregue as
          sequências extraídas e conceda recompensa adicional a um agente PPO
          durante o treinamento, quando este executar padrões de ação
          compatíveis.

    \item Treinar o agente PPO utilizando o Unity ML-Agents com o
          \textit{Reward Shaping} por proveniência e integrar o modelo
          resultante ao jogo em modo de inferência.

    \item Realizar uma avaliação quantitativa comparativa entre o chefe FSM e
          o chefe PPO, coletando métricas objetivas de desempenho em sessões
          controladas.
\end{enumerate}

\section{Justificativa}
\label{sec:justificativa}

A criação de NPCs com comportamentos interessantes, desafiadores e variados é
um problema recorrente no desenvolvimento de jogos. Embora técnicas
determinísticas como FSM continuem dominantes na indústria, a demanda por
experiências adaptativas motiva a investigação de abordagens baseadas em
aprendizado de máquina \cite{yannakakis2018artificial}.

O Aprendizado por Reforço oferece um arcabouço teórico robusto para esse fim,
mas sua aplicação prática em jogos esbarra em dois obstáculos frequentes: a
dificuldade de projetar funções de recompensa adequadas e a opacidade dos
comportamentos aprendidos. Este trabalho contribui para o enfrentamento de
ambos os obstáculos ao propor o uso de proveniência como ponte entre sessões
de jogo humanas e o treinamento de agentes, introduzindo um mecanismo que é,
ao mesmo tempo, automatizável --- pois não exige ajuste manual de recompensas
para cada cenário --- e interpretável --- pois as sequências que geram
recompensa adicional são derivadas de cadeias causais rastreáveis.

Além disso, a abordagem proposta preenche uma lacuna na literatura ao combinar
duas áreas tipicamente estudadas de forma independente: proveniência de dados
em jogos e aprendizado por reforço para NPCs. A integração entre essas áreas
permite transformar dados de \textit{gameplay} registrados com semântica
causal em sinais de treinamento, sem recorrer a demonstrações diretas ou
técnicas adversárias.

\section{Organização do Trabalho}
\label{sec:organizacao}

O restante deste trabalho está organizado da seguinte forma. O
\autoref{cap:referencial-teorico} apresenta o referencial teórico, abordando
os fundamentos de Máquinas de Estados Finitos, Aprendizado por Reforço, PPO,
Unity ML-Agents, Proveniência de Dados e \textit{Reward Shaping}. O
\autoref{cap:metodologia} descreve a metodologia adotada, incluindo o
pipeline de desenvolvimento, o protocolo de coleta de dados e as decisões de
escopo. O \autoref{cap:implementacao} detalha a implementação técnica do
\textit{vertical slice}, dos sistemas de proveniência e métricas, e da
integração do modelo treinado. O \autoref{cap:avaliacao} apresenta o cenário e
o protocolo de avaliação quantitativa. O \autoref{cap:resultados} expõe os
resultados obtidos na comparação entre os chefes FSM e PPO. O
\autoref{cap:discussao} discute os resultados, analisa as limitações do
trabalho e sugere direções para trabalhos futuros. Por fim, o
\autoref{cap:conclusao} sintetiza as contribuições do trabalho e apresenta as
considerações finais.

% =============================================================================
% Capítulo 2 — Referencial Teórico
% =============================================================================
\chapter{Referencial Teórico}
\label{cap:referencial-teorico}

Este capítulo apresenta os fundamentos teóricos que sustentam o presente
trabalho. São abordados os conceitos de Máquinas de Estados Finitos e sua
aplicação como técnica de controle de NPCs em jogos digitais, os princípios do
Aprendizado por Reforço com ênfase no algoritmo \textit{Proximal Policy
Optimization}, a plataforma Unity ML-Agents como ambiente de treinamento, o
conceito de Proveniência de Dados e sua aplicação ao rastreamento causal em
sessões de jogo, e a técnica de \textit{Reward Shaping} como mecanismo de
orientação do aprendizado.

% -----------------------------------------------------------------------------
\section{Máquinas de Estados Finitos em Jogos}
\label{sec:fsm}

Uma Máquina de Estados Finitos (FSM --- \textit{Finite State Machine}) é um
modelo computacional composto por um conjunto finito de estados, um conjunto de
transições entre esses estados e um estado inicial. Em cada instante, o sistema
encontra-se em exatamente um estado, e a ocorrência de um evento ou o
atendimento de uma condição pré-definida dispara uma transição para outro
estado \cite{millington2019ai}.

No contexto de jogos digitais, FSMs constituem uma das técnicas mais
difundidas para o controle de NPCs. Cada estado representa um comportamento
distinto do personagem --- como perseguir, atacar, recuar ou permanecer em
espera --- e as transições são determinadas por condições do ambiente, como a
distância até o jogador, o nível de vida remanescente ou a expiração de um
temporizador de \textit{cooldown}. A \autoref{fig:fsm-exemplo} ilustra, de
forma conceitual, a estrutura de uma FSM típica para um chefe de jogo 2D.

 \begin{figure}[htb]
     \centering
     \includegraphics[width=0.7\textwidth]{imagens/fsm_boss.png}
     \caption{Exemplo conceitual de FSM para um chefe de jogo 2D, com
              estados \textit{Idle}, \textit{Chase}, \textit{Attack} e
              \textit{Retreat}.}
     \label{fig:fsm-exemplo}
     \fonte{Elaborado pelo autor.}
 \end{figure}

As principais vantagens de FSMs residem em sua simplicidade de implementação,
previsibilidade comportamental e facilidade de depuração. Essas qualidades
tornam-nas particularmente adequadas para contextos de produção, nos quais os
projetistas necessitam de controle total sobre o comportamento do NPC e de
garantias quanto à reprodutibilidade de cenários de teste
\cite{millington2019ai}.

Contudo, a previsibilidade que constitui uma vantagem em termos de engenharia
converte-se em limitação sob a perspectiva do jogador: após poucas interações,
os padrões de comportamento tornam-se evidentes, reduzindo o desafio percebido
e comprometendo a rejogabilidade \cite{yannakakis2018artificial}. Além disso,
a complexidade de manutenção cresce de forma combinatória com o número de
estados e transições, tornando difícil a adição de novos comportamentos sem
introduzir efeitos colaterais indesejados.

Essas limitações motivam a investigação de abordagens complementares, capazes
de gerar comportamentos menos previsíveis e mais adaptáveis às ações do
jogador. Neste trabalho, a FSM é empregada como \textit{baseline}
determinístico, fornecendo um referencial estável e interpretável para
comparação com o agente treinado por Aprendizado por Reforço.

% -----------------------------------------------------------------------------
\section{Aprendizado por Reforço}
\label{sec:aprendizado-reforco}

O Aprendizado por Reforço (\textit{Reinforcement Learning} --- RL) é um
paradigma de aprendizado de máquina no qual um agente aprende a tomar decisões
por meio de interações com um ambiente. Diferentemente do aprendizado
supervisionado, no qual o modelo é treinado com pares de entrada e saída
rotulados, no RL o agente recebe sinais de recompensa numérica que indicam a
qualidade relativa de suas ações, sem instrução explícita sobre qual ação
tomar em cada situação \cite{sutton2018reinforcement}.

\subsection{Formulação como Processo de Decisão de Markov}
\label{subsec:mdp}

O arcabouço formal do RL é tipicamente modelado como um Processo de Decisão de
Markov (MDP --- \textit{Markov Decision Process}), definido pela tupla
$\langle S, A, P, R, \gamma \rangle$, em que:

\begin{itemize}
    \item $S$ é o conjunto de estados do ambiente;
    \item $A$ é o conjunto de ações disponíveis ao agente;
    \item $P(s' | s, a)$ é a função de transição, que define a probabilidade
          de o ambiente transitar para o estado $s'$ dado que o agente executou
          a ação $a$ no estado $s$;
    \item $R(s, a, s')$ é a função de recompensa, que atribui um valor
          numérico à transição de $s$ para $s'$ por meio da ação $a$;
    \item $\gamma \in [0, 1]$ é o fator de desconto, que pondera a importância
          relativa de recompensas futuras em relação a recompensas imediatas.
\end{itemize}

O objetivo do agente é aprender uma política $\pi(a|s)$ --- uma função que
mapeia estados a distribuições de probabilidade sobre ações --- que maximize o
retorno esperado acumulado ao longo do tempo. O retorno $G_t$ a partir do
instante $t$ é definido como:

\begin{equation}
    G_t = \sum_{k=0}^{\infty} \gamma^k \, R_{t+k+1}
    \label{eq:retorno}
\end{equation}

\noindent em que $R_{t+k+1}$ denota a recompensa recebida $k+1$ passos após o
instante $t$ \cite{sutton2018reinforcement}.

\subsection{Métodos Policy Gradient}
\label{subsec:policy-gradient}

Uma família importante de algoritmos de RL são os métodos de gradiente de
política (\textit{Policy Gradient}), nos quais a política $\pi_\theta$ é
parametrizada diretamente --- por exemplo, por uma rede neural com parâmetros
$\theta$ --- e otimizada por gradiente ascendente sobre uma função objetivo
que mede o retorno esperado. A vantagem dessa família de métodos é a
capacidade de lidar naturalmente com espaços de ação discretos e contínuos,
além de produzir políticas estocásticas que favorecem a exploração
\cite{sutton2018reinforcement}.

O gradiente da função objetivo, segundo o Teorema do Gradiente de Política, é
dado por:

\begin{equation}
    \nabla_\theta J(\theta)
    = \mathbb{E}_{\pi_\theta}
      \!\left[\,
        \nabla_\theta \log \pi_\theta(a_t | s_t) \; \hat{A}_t
      \,\right]
    \label{eq:policy-gradient}
\end{equation}

\noindent em que $\hat{A}_t$ é uma estimativa da função de vantagem
(\textit{advantage function}), que mede o quanto a ação $a_t$ é melhor ou pior
do que a ação média sob a política corrente.

\subsection{Proximal Policy Optimization (PPO)}
\label{subsec:ppo}

O \textit{Proximal Policy Optimization} (PPO), proposto por
\citeonline{schulman2017ppo}, é um algoritmo de gradiente de política que
endereça um problema recorrente em métodos anteriores: a instabilidade causada
por atualizações excessivamente grandes nos parâmetros da política. Algoritmos
como o \textit{Trust Region Policy Optimization} (TRPO) abordaram essa questão
por meio de restrições de segunda ordem, mas com custo computacional elevado.
O PPO propõe uma solução mais simples, baseada em uma função objetivo
\textit{clipped} que limita a magnitude das atualizações sem exigir cálculos
de segunda ordem.

A função objetivo do PPO é definida como:

\begin{equation}
    L^{CLIP}(\theta)
    = \mathbb{E}_t \!\left[\,
        \min\!\Big(
            r_t(\theta) \, \hat{A}_t, \;
            \text{clip}\big(r_t(\theta),\; 1 - \epsilon,\; 1 + \epsilon\big)
            \, \hat{A}_t
        \Big)
    \,\right]
    \label{eq:ppo-clip}
\end{equation}

\noindent em que $r_t(\theta) = \frac{\pi_\theta(a_t | s_t)}
{\pi_{\theta_{\text{old}}}(a_t | s_t)}$ é a razão entre a probabilidade da
ação sob a política atualizada e sob a política anterior, e
$\epsilon$ é um hiperparâmetro que controla o tamanho máximo permitido
para a atualização. Valores típicos de $\epsilon$ situam-se em torno de
$0{,}2$ \cite{schulman2017ppo}.

O mecanismo de \textit{clipping} opera da seguinte forma: quando a razão
$r_t(\theta)$ ultrapassa o intervalo $[1 - \epsilon, \; 1 + \epsilon]$, o
gradiente é atenuado, impedindo que a política se afaste excessivamente da
versão anterior em uma única iteração. Esse mecanismo garante atualizações
conservadoras, resultando em maior estabilidade durante o treinamento.

Do ponto de vista prático, o PPO apresenta características que o tornam
particularmente adequado para aplicações em jogos: estabilidade de treinamento
com ajuste mínimo de hiperparâmetros, compatibilidade com espaços de ação
discretos --- comuns em jogos 2D --- e eficiência amostral razoável em
ambientes simulados. Essas qualidades justificam sua adoção como algoritmo
padrão no Unity ML-Agents \cite{juliani2020unity}.

% -----------------------------------------------------------------------------
\section{Unity ML-Agents}
\label{sec:ml-agents}

O Unity ML-Agents \textit{Toolkit} é uma plataforma de código aberto
desenvolvida pela Unity Technologies que permite o treinamento de agentes
inteligentes em ambientes construídos com o motor de jogo Unity
\cite{juliani2020unity}. A plataforma integra o ambiente de simulação do Unity
com bibliotecas de aprendizado de máquina escritas em Python, notadamente o
PyTorch, estabelecendo um canal de comunicação bidirecional entre o motor de
jogo e o processo de treinamento.

\subsection{Arquitetura da Plataforma}
\label{subsec:ml-agents-arquitetura}

A arquitetura do ML-Agents organiza-se em três camadas principais:

\begin{enumerate}
    \item \textbf{Ambiente Unity}: contém a cena de simulação, os objetos de
          jogo e os scripts de comportamento dos agentes. Cada agente é
          implementado como um componente que herda da classe \texttt{Agent},
          sobrescrevendo métodos para coletar observações, executar ações e
          receber recompensas.

    \item \textbf{Comunicador}: protocolo de comunicação (\textit{side
          channel}) que transmite observações do Unity para o processo Python e
          retorna decisões de ação ao ambiente.

    \item \textbf{Treinador Python}: processo externo que executa o algoritmo
          de aprendizado (PPO, SAC, entre outros), atualiza os parâmetros da
          rede neural do agente e registra métricas de treinamento via
          TensorBoard.
\end{enumerate}

\subsection{Ciclo de Treinamento}
\label{subsec:ml-agents-ciclo}

O ciclo de treinamento no ML-Agents segue a dinâmica padrão do RL. A cada
passo de decisão (\textit{decision step}), o agente coleta observações do
ambiente --- como posições, velocidades e estados internos --- e as envia ao
treinador Python. O treinador seleciona uma ação de acordo com a política
corrente, e o agente a executa no ambiente Unity. O ambiente retorna uma
recompensa e o novo estado, fechando o ciclo.

Ao término do treinamento, o modelo resultante é exportado no formato ONNX
(\textit{Open Neural Network Exchange}) \cite{onnx2019}, um formato
interoperável que permite a execução de inferência diretamente no Unity, sem
a necessidade de manter o processo Python ativo. Essa característica é
fundamental para a integração do agente treinado em um jogo jogável, pois
permite que o modelo opere em modo de inferência local com desempenho
compatível com tempo real.

\subsection{Modos de Operação}
\label{subsec:ml-agents-modos}

O ML-Agents suporta três modos de operação relevantes para este trabalho:

\begin{itemize}
    \item \textbf{Treinamento}: o agente está conectado ao treinador Python e
          atualiza seus parâmetros a cada episódio;
    \item \textbf{Inferência} (\textit{InferenceOnly}): o agente utiliza um
          modelo ONNX previamente treinado para tomar decisões, sem conexão
          com o treinador;
    \item \textbf{Heurístico}: o agente é controlado por lógica programada
          manualmente, útil para coleta de demonstrações ou testes.
\end{itemize}

No presente trabalho, o treinamento foi realizado com PPO via treinador
Python, e o modelo final foi integrado ao Unity em modo
\textit{InferenceOnly} para a etapa de avaliação quantitativa.

% -----------------------------------------------------------------------------
\section{Proveniência de Dados}
\label{sec:proveniencia}

Proveniência de dados (\textit{Data Provenance}) refere-se ao registro
sistemático da origem, das transformações e do histórico de um dado ou
artefato ao longo de seu ciclo de vida. Em sua acepção mais geral, a
proveniência responde às perguntas ``de onde veio este dado?'' e ``por que ele
possui este valor?'', preservando as relações de causa e efeito entre eventos
e entidades em um fluxo de processamento \cite{provenance_base}.

\subsection{Proveniência versus Logs Convencionais}
\label{subsec:proveniencia-vs-logs}

É importante distinguir proveniência de dados de registros de log
convencionais. Enquanto um log tradicional armazena eventos sequenciais com
marcações temporais --- essencialmente respondendo ``o que aconteceu e
quando'' ---, um sistema de proveniência registra, adicionalmente, as relações
causais entre esses eventos. Cada evento possui referências explícitas ao
evento que o originou, formando uma estrutura de grafo dirigido acíclico (DAG)
que representa a cadeia causal da sessão.

Essa distinção é relevante porque a análise de logs planos permite apenas
correlações temporais (``o evento B ocorreu após o evento A''), ao passo que
a análise de proveniência permite inferências causais (``o evento B ocorreu
\textit{em consequência} do evento A''). Essa semântica adicional é
precisamente o que torna a proveniência útil como fonte de informação para
orientar processos de aprendizado.

\subsection{Proveniência em Jogos Digitais}
\label{subsec:proveniencia-jogos}

A aplicação de proveniência ao domínio de jogos digitais foi investigada por
\citeonline{provenance_base}, que propuseram o uso de grafos de proveniência
para a visualização e análise de sessões de jogo. Nessa abordagem, cada ação
relevante do \textit{gameplay} --- como ataques, movimentos, danos recebidos e
mortes --- é registrada como um nó do grafo, e as relações causais entre ações
e consequências são representadas como arestas direcionadas.

\citeonline{kohwalter2020provenance} ampliaram essa investigação com um
levantamento sobre aplicações de proveniência em jogos, demonstrando que os
grafos resultantes podem ser utilizados para análise \textit{post-mortem} de
sessões, identificação de padrões de \textit{gameplay}, detecção de
estratégias emergentes e depuração de \textit{game design}.

A estrutura típica de um grafo de proveniência em uma sessão de jogo de
combate segue um encadeamento causal como o ilustrado a seguir:

\begin{quote}
    \texttt{SessionStart} $\rightarrow$ \texttt{BossAttack} $\rightarrow$
    \texttt{PlayerDamageTaken} $\rightarrow$ \texttt{PlayerAttack}
    $\rightarrow$ \texttt{BossDamageTaken} $\rightarrow$
    \texttt{BossDeath} $\rightarrow$ \texttt{SessionEnd}
\end{quote}

Nessa cadeia, cada evento possui uma referência ao evento que o precedeu
causalmente. A relação entre \texttt{PlayerAttack} e
\texttt{BossDamageTaken}, por exemplo, não é meramente temporal: ela indica
que o dano no chefe foi \textit{causado} pelo ataque do jogador. Essa
distinção é fundamental para a extração de padrões comportamentais
significativos.

No presente trabalho, a proveniência é empregada em dois momentos distintos:
como dado bruto exportado em formato JSON para análise e documentação das
sessões de jogo; e como fonte indireta de informação para o treinamento do
agente, por meio da extração de sequências de ações associadas a desfechos
vitoriosos.

% -----------------------------------------------------------------------------
\section{Modelagem de Recompensa (\textit{Reward Shaping})}
\label{sec:reward-shaping}

A definição de funções de recompensa é uma das etapas mais delicadas do
treinamento por Aprendizado por Reforço. Uma função de recompensa mal
projetada pode levar a comportamentos degenerados: o agente pode encontrar
formas de maximizar a recompensa que não correspondem ao objetivo desejado
pelo projetista --- fenômeno conhecido como \textit{reward hacking}
\cite{sutton2018reinforcement}.

\subsection{Conceito e Fundamentação Teórica}
\label{subsec:reward-shaping-conceito}

\textit{Reward Shaping} é uma técnica que consiste na adição de uma
recompensa auxiliar $F(s, a, s')$ à recompensa original do ambiente $R(s, a,
s')$, com o objetivo de fornecer sinais intermediários que orientem o agente
em direção ao comportamento desejado sem alterar a política ótima do problema
original \cite{ng1999reward}.

\citeonline{ng1999reward} demonstraram formalmente que, quando a recompensa
auxiliar $F$ assume a forma de uma função potencial baseada em estados:

\begin{equation}
    F(s, a, s') = \gamma \, \Phi(s') - \Phi(s)
    \label{eq:pbrs}
\end{equation}

\noindent em que $\Phi: S \rightarrow \mathbb{R}$ é uma função potencial e
$\gamma$ é o fator de desconto, a política ótima do MDP original é preservada.
Essa propriedade, conhecida como invariância de política sob \textit{Potential-
Based Reward Shaping} (PBRS), garante que o \textit{Reward Shaping} não
introduz vícios permanentes no aprendizado: ele apenas acelera a convergência
em direção à política ótima, sem desviá-la.

\subsection{Reward Shaping versus Imitação Direta}
\label{subsec:shaping-vs-imitacao}

É relevante distinguir \textit{Reward Shaping} de técnicas de aprendizagem
por imitação, como GAIL \cite{ho2016gail}. Na imitação direta, o agente
aprende a reproduzir trajetórias demonstradas por um especialista, utilizando
um discriminador adversário que diferencia o comportamento do agente do
comportamento demonstrado. Embora eficaz, essa abordagem vincula o agente ao
comportamento observado, limitando sua capacidade de exploração e dificultando
a emergência de estratégias originais.

O \textit{Reward Shaping}, por outro lado, atua como um sinal orientador que
não substitui a recompensa do ambiente. O agente continua explorando
livremente e otimizando a recompensa total; a componente adicional apenas
indica que certos padrões de comportamento são preferíveis, sem impor uma
trajetória específica. Essa distinção é particularmente relevante no contexto
do presente trabalho, no qual a intenção não é que o chefe \textit{copie} o
comportamento de um jogador humano, mas sim que receba orientação de padrões
causais extraídos de sessões bem-sucedidas.

\subsection{Aplicação ao Contexto do Trabalho}
\label{subsec:shaping-aplicacao}

No pipeline proposto neste trabalho, o \textit{Reward Shaping} é alimentado
por sequências de ações extraídas de grafos de proveniência. Essas sequências
representam padrões de comportamento que precederam eventos favoráveis ---
especificamente, dano causado no chefe em sessões vitoriosas. Quando o agente
PPO, durante o treinamento, executa uma sequência de ações que coincide com
uma das sequências extraídas, ele recebe uma recompensa adicional.

Esse mecanismo combina as vantagens das duas áreas: a proveniência fornece
informação causal interpretável sobre o que constituiu comportamento bem-
sucedido, e o \textit{Reward Shaping} converte essa informação em um sinal
de treinamento compatível com o framework de RL, sem sacrificar a autonomia
exploratória do agente.

% -----------------------------------------------------------------------------
\section{Considerações sobre o Referencial}
\label{sec:consideracoes-referencial}

Os conceitos apresentados neste capítulo fornecem a base teórica para a
compreensão do pipeline proposto. A FSM oferece o \textit{baseline}
determinístico necessário para uma comparação objetiva. O Aprendizado por
Reforço, por meio do PPO, fornece o algoritmo de treinamento do agente. O
Unity ML-Agents viabiliza a execução desse treinamento diretamente no motor
de jogo. A Proveniência de Dados possibilita o registro e a análise causal
de sessões de jogo. E o \textit{Reward Shaping} fornece o mecanismo teórico
para converter padrões causais em sinais de recompensa auxiliares, preservando
a política ótima.

A integração dessas cinco áreas constitui o diferencial metodológico do
presente trabalho, cujos detalhes de implementação e aplicação são
apresentados nos capítulos subsequentes.

% =============================================================================
% Capítulo 3 — Metodologia
% =============================================================================
\chapter{Metodologia}
\label{cap:metodologia}

Este capítulo descreve as escolhas metodológicas que orientaram o
desenvolvimento do presente trabalho. São apresentados a concepção do
protótipo jogável, o ambiente técnico e as ferramentas utilizadas, o
pipeline completo de coleta, extração e treinamento, o protocolo de avaliação
quantitativa e, por fim, as decisões de escopo que delimitaram a contribuição
do trabalho.

% -----------------------------------------------------------------------------
\section{Concepção do Protótipo (\textit{Vertical Slice})}
\label{sec:vertical-slice}

A estratégia de desenvolvimento adotada foi a construção de um
\textit{vertical slice}: um recorte funcional do jogo que implementa, de
ponta a ponta, todas as camadas necessárias para validar a proposta --- desde
a jogabilidade até a exportação de métricas --- sem a pretensão de representar
um produto comercial completo. Essa abordagem é recorrente na indústria de
jogos quando o objetivo é demonstrar a viabilidade de um conceito em escopo
controlado \cite{yannakakis2018artificial}.

O protótipo, denominado \textit{Chronicles Of The Lost Word}, consiste em um
jogo 2D de ação construído no motor Unity. O recorte implementado concentra-se
exclusivamente em uma arena de combate contra um chefe (\textit{boss}),
excluindo elementos como exploração de fases, sistemas de progressão,
narrativa interativa ou interfaces de menu. Essa restrição é intencional: o
escopo reduzido permite isolar a variável de interesse --- o comportamento do
chefe --- e avaliar o impacto do pipeline de proveniência e \textit{Reward
Shaping} sem interferência de sistemas periféricos.

\subsection{O Personagem Jogável (Elian)}
\label{subsec:elian}

O personagem controlado pelo jogador, Elian, possui um conjunto de mecânicas
de ação que define o espaço de possibilidades da interação com o chefe:

\begin{itemize}
    \item \textbf{Movimento horizontal}: deslocamento à esquerda e à direita
          no plano da arena;
    \item \textbf{Pulo} (\textit{PlayerJump}): salto vertical com detecção
          de solo via \textit{ground check};
    \item \textbf{Ataque} (\textit{PlayerAttack}): golpe corpo a corpo com
          área de efeito definida por \texttt{OverlapBoxAll}, aplicando dano
          em alvos que implementam a interface \texttt{IDamageable};
    \item \textbf{Dash} (\textit{PlayerDash}): avanço rápido horizontal com
          breve período de invulnerabilidade;
    \item \textbf{Vida} (\textit{PlayerHealth}): sistema de pontos de vida
          com dano recebido, morte e eventos associados.
\end{itemize}

Essas três ações ofensivas e evasivas --- pulo, ataque e \textit{dash} ---
constituem o vocabulário de ações registrado pelo sistema de proveniência e,
posteriormente, utilizado na extração de sequências vitoriosas. A escolha
desse conjunto restrito foi deliberada: um vocabulário de ação compacto
facilita a identificação de padrões causais nos grafos de proveniência e
mantém o espaço de sequências possíveis tratável computacionalmente.

\subsection{A Arena de Combate}
\label{subsec:arena}

A arena consiste em um espaço horizontal delimitado por paredes invisíveis,
com plataforma plana e sem obstáculos adicionais. Esse \textit{design}
minimalista garante que as diferenças de comportamento entre os chefes FSM e
PPO reflitam exclusivamente suas estratégias de combate, sem influência de
variáveis ambientais como plataformas elevadas, armadilhas ou itens coletáveis.

Cada sessão de combate é gerenciada pelo componente \texttt{ArenaManager}, que
controla o ciclo de vida da sessão: início, detecção de vitória (morte do
chefe) ou derrota (morte do jogador), e encerramento com exportação automática
de métricas.

% -----------------------------------------------------------------------------
\section{Arquitetura e Ferramentas}
\label{sec:arquitetura-ferramentas}

O desenvolvimento do protótipo apoiou-se em um conjunto de ferramentas
selecionadas por sua adequação ao domínio de jogos 2D com aprendizado por
reforço. A \autoref{tab:ambiente-tecnico} sintetiza as versões utilizadas.

\begin{table}[htb]
    \centering
    \caption{Ambiente técnico utilizado no desenvolvimento e treinamento.}
    \label{tab:ambiente-tecnico}
    \begin{tabular}{l l}
        \toprule
        \textbf{Componente}         & \textbf{Versão / Identificação} \\
        \midrule
        Unity Engine                & 6000.3.17f1                     \\
        Linguagem de scripts        & C\#                             \\
        ML-Agents (pacote Unity)    & 4.0.3                           \\
        Input System (pacote Unity) & 1.19.0                          \\
        Test Framework (Unity)      & 1.6.0                           \\
        Python                      & 3.10.12                         \\
        \texttt{mlagents} (Python)  & 1.1.0                           \\
        PyTorch                     & 2.8.0+cpu                       \\
        ONNX                        & 1.15.0                          \\
        \bottomrule
    \end{tabular}
    \fonte{Elaborado pelo autor.}
\end{table}

\subsection{Unity como Motor de Jogo e Ambiente de Simulação}
\label{subsec:unity-motor}

O Unity \cite{unity2024} foi escolhido como motor de jogo por reunir
capacidades de desenvolvimento 2D nativo --- incluindo sistemas de física
(\texttt{Rigidbody2D}, \texttt{BoxCollider2D}), renderização por
\textit{sprites} e \textit{Input System} moderno --- com suporte integrado ao
\textit{toolkit} de aprendizado de máquina ML-Agents
\cite{juliani2020unity}. Essa integração nativa elimina a necessidade de
adaptadores externos entre o ambiente de simulação e o processo de
treinamento, simplificando o pipeline.

Toda a lógica do jogo foi implementada em C\#, organizada em módulos com
responsabilidades claras: \textit{Player}, \textit{Boss FSM}, \textit{Boss
PPO}, Proveniência, Arena e Métricas. Cada módulo comunica-se por meio de
eventos (\textit{delegates} e \texttt{Action<T>} do C\#), seguindo o
princípio de comunicação dirigida por eventos para minimizar o acoplamento
entre componentes.

\subsection{Unity ML-Agents como Ponte para RL}
\label{subsec:ml-agents-ponte}

O Unity ML-Agents \textit{Toolkit} \cite{juliani2020unity} atua como
intermediário entre o ambiente de jogo e os algoritmos de aprendizado por
reforço. No lado Unity, o agente é implementado como um componente C\# que
herda da classe \texttt{Agent}, expondo métodos para coleta de observações
(\texttt{CollectObservations}), execução de ações
(\texttt{OnActionReceived}) e atribuição de recompensas
(\texttt{AddReward}). No lado Python, o treinador executa o algoritmo PPO
\cite{schulman2017ppo} utilizando PyTorch como \textit{backend} de
computação, comunicando-se com o Unity via protocolo gRPC.

Ao término do treinamento, o modelo é exportado em formato ONNX
\cite{onnx2019}, permitindo inferência diretamente no Unity sem dependência
do processo Python. Essa separação entre fase de treinamento e fase de
inferência é essencial para a avaliação: o modelo avaliado opera de forma
autônoma, determinística e reprodutível.

% -----------------------------------------------------------------------------
\section{Pipeline de Treinamento e Coleta}
\label{sec:pipeline}

O pipeline de desenvolvimento segue um fluxo sequencial de seis etapas, cada
uma produzindo artefatos consumidos pela etapa seguinte. A
\autoref{fig:pipeline} ilustra esse fluxo de forma esquemática.

 \begin{figure}[htb]
     \centering
     \includegraphics[width=0.85\textwidth]{imagens/pipeline-treinamento}
     \caption{Pipeline de treinamento e coleta: do \textit{gameplay} humano
              à avaliação quantitativa FSM vs.\ PPO.}
     \label{fig:pipeline}
     \fonte{Elaborado pelo autor.}
 \end{figure}

\subsection{Etapa 1 --- Coleta de Sessões contra o Chefe FSM}
\label{subsec:etapa1}

Na primeira etapa, o desenvolvedor joga múltiplas sessões como Elian contra o
chefe controlado por FSM. O chefe FSM opera com quatro estados determinísticos
--- \textit{Idle}, \textit{Chase}, \textit{Attack} e \textit{Retreat} ---
com parâmetros fixos de velocidade, alcance de ataque e \textit{cooldown}. A
previsibilidade do FSM é uma característica desejável nesta etapa, pois
garante que as sessões coletadas reflitam o comportamento do jogador humano
em condições estáveis e reprodutíveis.

\subsection{Etapa 2 --- Registro de Proveniência}
\label{subsec:etapa2}

Durante cada sessão, o sistema de proveniência registra automaticamente os
eventos de \textit{gameplay} como nós de um grafo causal. Cada nó contém
atributos padronizados: identificador do evento (\texttt{eventId}), marcação
temporal (\texttt{timestamp}), identificador do ator (\texttt{actorId}), tipo
de ação (\texttt{actionType}), posição no espaço (\texttt{position}), valor
associado (\texttt{value}) e referência ao evento causador
(\texttt{parentEventId}). Essa última referência é a ligação causal que
transforma a sequência temporal de eventos em uma cadeia causal explícita.

O componente \texttt{ProvenanceLogger} escuta eventos emitidos pelos scripts
de \textit{gameplay} --- \texttt{PlayerController}, \texttt{PlayerCombat},
\texttt{PlayerHealth}, \texttt{BossHealth} e \texttt{BossFsmController} ---
e registra cada um no grafo em memória, gerenciado pelo componente
\texttt{ProvenanceGraph}. Ao término da sessão, o
\texttt{ProvenanceExporter} serializa o grafo completo em formato JSON.

\subsection{Etapa 3 --- Extração de Sequências Vitoriosas}
\label{subsec:etapa3}

A extração de padrões comportamentais a partir dos grafos de proveniência é
realizada por um script Python dedicado. Esse script opera em três passos
sequenciais:

\begin{enumerate}
    \item \textbf{Filtragem}: lê os arquivos JSON exportados e seleciona
          apenas as sessões cujo campo \texttt{result} é igual a
          \texttt{"victory"}, descartando sessões encerradas com derrota ou
          incompletas;
    \item \textbf{Localização de eventos-alvo}: percorre os nós de cada
          grafo filtrado e identifica todas as ocorrências do evento
          \texttt{BossDamageTaken}, que indica que o chefe recebeu dano;
    \item \textbf{Extração de sequências}: para cada evento
          \texttt{BossDamageTaken}, retrocede na cadeia causal e coleta as
          três ações imediatamente anteriores pertencentes ao vocabulário de
          ações permitidas (\texttt{PlayerJump}, \texttt{PlayerAttack},
          \texttt{PlayerDash}), formando uma sequência de tamanho fixo.
\end{enumerate}

A \autoref{tab:config-extracao} resume os parâmetros utilizados nesse
processo.

\begin{table}[htb]
    \centering
    \caption{Configuração do processo de extração de sequências vitoriosas.}
    \label{tab:config-extracao}
    \begin{tabular}{l l}
        \toprule
        \textbf{Parâmetro}    & \textbf{Valor}                                  \\
        \midrule
        Ações permitidas      & \texttt{PlayerJump}, \texttt{PlayerAttack},
                                \texttt{PlayerDash}                              \\
        Tamanho da sequência  & 3                                                \\
        Evento-alvo           & \texttt{BossDamageTaken}                         \\
        Formato de saída      & JSON (\texttt{winning\_sequences.json})          \\
        \bottomrule
    \end{tabular}
    \fonte{Elaborado pelo autor.}
\end{table}

O resultado é um arquivo JSON contendo a lista de sequências de ações que, em
sessões vitoriosas, precederam causalmente eventos de dano no chefe. Esse
artefato não constitui uma demonstração direta de \textit{gameplay}: ele
codifica \textit{padrões} de ação associados a desfechos favoráveis, sem
registrar trajetórias completas de estados e ações.

\subsection{Etapa 4 --- Treinamento PPO com Reward Shaping por Proveniência}
\label{subsec:etapa4}

O agente PPO foi treinado utilizando o Unity ML-Agents com o algoritmo PPO
como treinador principal \cite{schulman2017ppo}. A configuração do
treinamento foi definida em um arquivo YAML consumido pelo treinador Python.
A \autoref{tab:config-ppo} apresenta os hiperparâmetros utilizados.

\begin{table}[htb]
    \centering
    \caption{Hiperparâmetros do treinamento PPO.}
    \label{tab:config-ppo}
    \begin{tabular}{l r}
        \toprule
        \textbf{Hiperparâmetro}          & \textbf{Valor}  \\
        \midrule
        \texttt{trainer\_type}           & PPO             \\
        \texttt{batch\_size}             & 1\,024          \\
        \texttt{buffer\_size}            & 10\,240         \\
        \texttt{learning\_rate}          & $3{,}0 \times 10^{-4}$ \\
        \texttt{beta} (entropia)         & $5{,}0 \times 10^{-3}$ \\
        \texttt{epsilon} (clipping)      & 0,2             \\
        \texttt{lambd} (GAE-$\lambda$)   & 0,95            \\
        \texttt{num\_epoch}              & 3               \\
        \texttt{hidden\_units}           & 128             \\
        \texttt{num\_layers}             & 2               \\
        \texttt{max\_steps}              & 500\,000        \\
        \texttt{time\_horizon}           & 64              \\
        Sinal de recompensa              & Extrínseco      \\
        \bottomrule
    \end{tabular}
    \fonte{Elaborado pelo autor.}
\end{table}

Durante o treinamento, o componente \texttt{ProvenanceRewardShaper} mantém um
\textit{buffer} circular das últimas ações executadas pelo agente. A cada
passo de decisão, o \textit{buffer} é comparado com as sequências vitoriosas
carregadas do arquivo JSON. Quando há correspondência exata entre o
\textit{buffer} e alguma sequência vitoriosa, uma recompensa adicional de
$+2{,}0$ é concedida ao agente via \texttt{AddReward}. Essa recompensa
auxiliar opera em conjunto com a recompensa extrínseca do ambiente, sem
substituí-la.

Um aspecto de projeto relevante é a correspondência de vocabulário entre as
ações do agente PPO e as ações humanas registradas na proveniência. As ações
discretas do agente --- pulo, ataque e \textit{dash} --- registram no
\textit{Reward Shaper} as mesmas \textit{strings} utilizadas pelo sistema de
proveniência (\texttt{PlayerJump}, \texttt{PlayerAttack},
\texttt{PlayerDash}), garantindo compatibilidade direta na comparação de
sequências.

O \textit{run} final de treinamento, identificado como
\texttt{boss\_v1\_500k}, alcançou 500.000 passos de treinamento e produziu o
modelo \texttt{BossAgent.onnx}, que foi exportado e integrado ao projeto
Unity em modo \texttt{InferenceOnly} para a etapa de avaliação.

\subsection{Etapa 5 --- Integração e Preparação para Avaliação}
\label{subsec:etapa5}

O modelo ONNX resultante foi integrado ao projeto Unity e associado ao
\textit{prefab} do chefe PPO. Antes da coleta oficial de dados, foi
necessário ajustar as cenas de avaliação para garantir condições equivalentes
entre as duas condições experimentais:

\begin{itemize}
    \item Adição do componente \texttt{PlayerHealth} com 100 pontos de vida
          ao personagem Elian, permitindo que sessões terminassem por derrota
          do jogador;
    \item Implementação do componente \texttt{BossAttackDamage}, que aplica
          dano real de 10 pontos quando o chefe (FSM ou PPO) ataca dentro do
          alcance, com \textit{cooldown} de 0,8 segundos;
    \item Equalização de parâmetros de dano e alcance entre as cenas FSM e
          PPO, eliminando vieses de configuração.
\end{itemize}

\subsection{Etapa 6 --- Avaliação Quantitativa}
\label{subsec:etapa6}

A etapa final consiste na coleta controlada de sessões jogáveis para
comparação. O protocolo de avaliação define a realização de 10 sessões contra
o chefe FSM e 10 sessões contra o chefe PPO, todas jogadas pelo
desenvolvedor em condições padronizadas. As métricas são exportadas
automaticamente em formato CSV pelo sistema \texttt{GameMetrics} e
\texttt{GameMetricsExporter}, incluindo duração da sessão, dano recebido
pelo chefe, dano recebido pelo jogador, contagem de ações do jogador e
contagem de ações do chefe.

A análise dos dados é apresentada nos capítulos de Avaliação e Resultados.

% -----------------------------------------------------------------------------
\section{Protocolo de Avaliação}
\label{sec:protocolo-avaliacao}

A avaliação quantitativa segue um protocolo controlado, projetado para
maximizar a comparabilidade entre as duas condições experimentais dentro das
restrições de escopo do trabalho.

\subsection{Condições Experimentais}
\label{subsec:condicoes}

Duas condições foram definidas:

\begin{itemize}
    \item \textbf{Condição FSM}: o jogador enfrenta o chefe controlado por
          Máquina de Estados Finitos, na cena \texttt{ArenaChefe.unity};
    \item \textbf{Condição PPO}: o jogador enfrenta o chefe controlado pelo
          modelo PPO treinado, em modo \texttt{InferenceOnly}, na cena
          \texttt{BossArena\_PPO.unity}.
\end{itemize}

Em ambas as condições, o personagem Elian é controlado pelo desenvolvedor com
as mesmas mecânicas e parâmetros. Os parâmetros de dano do chefe (10 pontos
por ataque, alcance de 1,25 unidades, \textit{cooldown} de 0,8 segundos)
e a vida do jogador (100 pontos) são idênticos nas duas cenas, garantindo
que diferenças observadas nas métricas reflitam exclusivamente o
comportamento do chefe.

\subsection{Tamanho da Amostra}
\label{subsec:amostra}

A amostra oficial é composta por 20 sessões: 10 na condição FSM e 10 na
condição PPO. Embora o tamanho da amostra seja limitado, ele é suficiente para
identificar diferenças comportamentais mensuráveis entre as duas condições,
que é o objetivo do trabalho. Não se pretende, com essa amostra, estabelecer
significância estatística robusta ou generalizar resultados para outros
contextos --- essas restrições são discutidas como limitações no
\autoref{cap:discussao}.

\subsection{Métricas Coletadas}
\label{subsec:metricas}

As métricas são registradas automaticamente pelo sistema
\texttt{GameMetrics} e exportadas em CSV com formato invariante. As métricas
principais são:

\begin{itemize}
    \item \textbf{Resultado da sessão}: vitória, derrota ou encerramento
          prematuro;
    \item \textbf{Duração da sessão} (em segundos);
    \item \textbf{Dano total recebido pelo chefe}
          (\texttt{boss\_damage\_taken});
    \item \textbf{Dano total recebido pelo jogador}
          (\texttt{player\_damage\_taken});
    \item \textbf{Contagem de ações do jogador}: pulos, ataques e
          \textit{dashes};
    \item \textbf{Contagem de ações do chefe}: ataques realizados; para o
          PPO, inclui-se também a distribuição completa das seis ações
          discretas (Idle, Move Left, Move Right, Jump, Attack, Dash).
\end{itemize}

% -----------------------------------------------------------------------------
\section{Decisões de Escopo e Limitações Metodológicas}
\label{sec:decisoes-escopo}

O escopo do trabalho foi delimitado por decisões deliberadas, tomadas ao
longo do desenvolvimento para manter a viabilidade do TCC sem comprometer a
integridade dos resultados. Esta seção documenta as decisões mais relevantes
e suas justificativas.

\subsection{Remoção da Aprendizagem por Imitação (GAIL)}
\label{subsec:remocao-gail}

O plano inicial do trabalho previa a utilização de \textit{Generative
Adversarial Imitation Learning} (GAIL) \cite{ho2016gail} como mecanismo
complementar de treinamento, a partir de demonstrações gravadas em formato
\texttt{.demo} do ML-Agents. Essa abordagem foi removida do escopo por duas
razões:

\begin{enumerate}
    \item O GAIL introduziria uma dependência de demonstrações diretas,
          conflitando com a proposta central de utilizar proveniência --- e
          não imitação --- como fonte de orientação para o aprendizado;
    \item A coexistência de GAIL e \textit{Reward Shaping} por proveniência
          dificultaria a atribuição de efeitos: seria impossível determinar
          se o comportamento aprendido resultou da imitação, do
          \textit{Reward Shaping} ou de sua combinação.
\end{enumerate}

A remoção do GAIL reforça a contribuição acadêmica do trabalho ao isolar o
efeito da proveniência como único mecanismo auxiliar de orientação.

\subsection{Remoção de Playtesters Externos}
\label{subsec:remocao-playtesters}

O plano original previa a participação de jogadores externos
(\textit{playtesters}) e a aplicação de questionários de experiência. Essa
etapa foi removida por duas razões de escopo:

\begin{enumerate}
    \item A logística de recrutamento, consentimento e coleta presencial de
          dados ultrapassava o prazo disponível para o TCC;
    \item A inclusão de variáveis subjetivas (percepção de dificuldade,
          satisfação) ampliaria o escopo analítico para além do que o
          trabalho se propõe a endereçar.
\end{enumerate}

Em substituição, a avaliação foi redesenhada como uma coleta quantitativa
controlada, conduzida pelo próprio desenvolvedor. Essa escolha é uma limitação
reconhecida --- o avaliador conhece os padrões de ambos os chefes --- mas
garante padronização absoluta das condições de teste e viabilidade temporal.

\subsection{Modelo Congelado na Avaliação}
\label{subsec:modelo-congelado}

O modelo PPO utilizado na avaliação é o resultado do \textit{run} final de
treinamento (\texttt{boss\_v1\_500k}), exportado como arquivo ONNX e
integrado ao Unity em modo \texttt{InferenceOnly}. Após a exportação, o
modelo não foi retreinado, ajustado ou modificado em nenhuma circunstância,
inclusive após a adição do sistema de vida real do jogador
(\texttt{PlayerHealth}).

Essa decisão foi deliberada: a avaliação deve medir o comportamento do modelo
congelado, não de um modelo ajustado \textit{post hoc} com base nos
resultados observados. Como consequência, o modelo avaliado não observa a
vida real do jogador como entrada --- essa observação era um
\textit{placeholder} na arquitetura utilizada durante o treino. Essa limitação
é documentada e discutida no \autoref{cap:discussao}.

\subsection{Preservação do Comportamento Repetitivo}
\label{subsec:preservacao-spam}

Durante as sessões de avaliação, observou-se que o agente PPO apresenta alta
frequência de ações de ataque (\textit{spam}), resultando em comportamento
visualmente repetitivo. Embora uma intervenção no código pudesse limitar essa
frequência --- por exemplo, adicionando um \textit{cooldown} forçado às
ações --- optou-se por preservar o comportamento original do modelo.

A justificativa é metodológica: o comportamento repetitivo é um resultado
experimental do modelo congelado. Corrigi-lo antes da análise
descaracterizaria a avaliação, introduzindo uma modificação \textit{post hoc}
que alteraria os dados coletados sem refletir uma melhoria genuína no
aprendizado do agente.

\subsection{Síntese das Decisões}
\label{subsec:sintese-decisoes}

A \autoref{tab:decisoes-escopo} sintetiza as decisões de escopo e suas
justificativas.

\begin{table}[htb]
    \centering
    \caption{Decisões de escopo e justificativas.}
    \label{tab:decisoes-escopo}
    \begin{tabular}{p{5.5cm} p{8cm}}
        \toprule
        \textbf{Decisão} & \textbf{Justificativa} \\
        \midrule
        Remover GAIL e demonstrações &
            Preservar a proveniência como único sinal auxiliar, isolando
            a contribuição acadêmica. \\
        \addlinespace
        Remover playtesters e questionários &
            Viabilidade temporal e foco em avaliação quantitativa objetiva. \\
        \addlinespace
        Não retreinar o modelo após adicionar vida real &
            Garantir que a avaliação meça o modelo congelado, sem
            ajustes \textit{post hoc}. \\
        \addlinespace
        Preservar o \textit{spam} de ataque do PPO &
            O comportamento é resultado experimental; corrigi-lo
            descaracterizaria a avaliação. \\
        \addlinespace
        Desabilitar JSON de proveniência na cena PPO &
            A métrica principal é o CSV; o JSON da cena FSM já
            evidencia o mecanismo de proveniência. \\
        \bottomrule
    \end{tabular}
    \fonte{Elaborado pelo autor.}
\end{table}

Essas decisões refletem a prioridade de manter a integridade metodológica do
trabalho dentro das restrições práticas de um Trabalho de Conclusão de Curso.
As limitações decorrentes são reconhecidas e discutidas no
\autoref{cap:discussao}, juntamente com sugestões de trabalhos futuros que
poderiam endereçá-las.

% =============================================================================
% Capítulo 4 — Implementação
% =============================================================================
\chapter{Implementação}
\label{cap:implementacao}

Este capítulo descreve os detalhes técnicos da implementação do
\textit{vertical slice}, apresentando a organização do código em módulos, os
sistemas de \textit{gameplay} que compõem o \textit{baseline} e o personagem
jogável, o núcleo de proveniência responsável pelo registro causal das
sessões, a integração com aprendizado por reforço via \textit{Reward Shaping}
e \texttt{BossAgent}, e os sistemas de coleta de métricas e gerenciamento de
sessões.

% -----------------------------------------------------------------------------
\section{Arquitetura do Código}
\label{sec:arquitetura-codigo}

A base de código foi organizada em módulos com responsabilidades claramente
delimitadas, seguindo o princípio de separação de responsabilidades
(\textit{Separation of Concerns}). Cada módulo encapsula uma área funcional
do sistema e expõe sua interface pública por meio de eventos C\#
(\texttt{Action<T>} e \textit{delegates}), minimizando o acoplamento direto
entre componentes.

A \autoref{tab:modulos} apresenta os módulos e suas responsabilidades.

\begin{table}[htb]
    \centering
    \caption{Módulos da arquitetura e suas responsabilidades.}
    \label{tab:modulos}
    \begin{tabular}{l p{9.5cm}}
        \toprule
        \textbf{Módulo}     & \textbf{Responsabilidade} \\
        \midrule
        Player              & Captura de entrada, movimento, ataque, \textit{dash}
                              e sistema de vida do personagem Elian. \\
        \addlinespace
        Boss FSM             & \textit{Baseline} determinístico com estados
                              \textit{Idle}, \textit{Chase}, \textit{Attack} e
                              \textit{Retreat}. \\
        \addlinespace
        Proveniência         & Registro causal de eventos em memória, exportação
                              para JSON e módulo de \textit{Reward Shaping}. \\
        \addlinespace
        Boss PPO             & Agente ML-Agents com observações vetoriais, ações
                              discretas e integração com o
                              \texttt{ProvenanceRewardShaper}. \\
        \addlinespace
        Arena / Métricas     & Ciclo de vida da sessão (início, vitória, derrota),
                              coleta de métricas e exportação CSV. \\
        \bottomrule
    \end{tabular}
    \fonte{Elaborado pelo autor.}
\end{table}

\subsection{Comunicação por Eventos}
\label{subsec:comunicacao-eventos}

A comunicação entre módulos é inteiramente dirigida por eventos. Cada
componente que produz informação relevante --- como a execução de um ataque,
a recepção de dano ou a morte de um personagem --- expõe um evento público
(\texttt{Action} ou \texttt{Action<T>}) ao qual outros componentes podem se
inscrever. Esse padrão arquitetural garante que:

\begin{itemize}
    \item O emissor do evento desconhece seus consumidores, eliminando
          dependências circulares;
    \item Novos sistemas (como o de proveniência ou o de métricas) podem ser
          adicionados sem modificar os scripts de \textit{gameplay} existentes,
          bastando inscrever-se nos eventos já expostos;
    \item A testabilidade dos componentes é preservada, pois cada módulo pode
          ser verificado isoladamente por meio de testes unitários que disparam
          ou escutam eventos sem depender da cena completa do Unity.
\end{itemize}

Por exemplo, o evento \texttt{OnBossAttackPerformed}, emitido pelo
\texttt{BossFsmController} quando o chefe executa um ataque, é consumido
simultaneamente pelo \texttt{ProvenanceLogger} (para registrar o evento no
grafo), pelo \texttt{GameMetrics} (para incrementar a contagem de ataques) e
pelo \texttt{BossAttackDamage} (para aplicar dano ao jogador). Nenhum desses
consumidores conhece os demais.

% -----------------------------------------------------------------------------
\section{Sistemas de Gameplay e Baseline}
\label{sec:sistemas-gameplay}

Esta seção descreve a implementação dos scripts que compõem o personagem
jogável e o chefe \textit{baseline} controlado por FSM.

\subsection{Personagem Jogável --- PlayerController}
\label{subsec:player-controller}

O \texttt{PlayerController} é o componente central do personagem Elian. Ele
captura as entradas do jogador por meio do \textit{Unity Input System} e
traduz essas entradas em ações físicas executadas via \texttt{Rigidbody2D}.
A separação entre captura de \textit{input} e execução de física segue as
diretrizes de arquitetura do projeto: as entradas são lidas no \texttt{Update}
e as forças são aplicadas no \texttt{FixedUpdate}, garantindo consistência
com o \textit{timestep} fixo da simulação física.

As ações implementadas são:

\begin{itemize}
    \item \textbf{Movimento horizontal}: aplica velocidade ao
          \texttt{Rigidbody2D} na direção indicada pelo eixo de entrada;
    \item \textbf{Pulo}: aplica impulso vertical quando o personagem está em
          contato com o solo, verificado por um \textit{ground check}
          via \textit{raycast} ou \textit{overlap};
    \item \textbf{Dash}: desloca o personagem horizontalmente com velocidade
          aumentada por um curto período;
    \item \textbf{Ataque}: delega a lógica de dano ao componente
          \texttt{PlayerCombat}.
\end{itemize}

Cada ação emite um evento público quando executada --- por exemplo,
\texttt{OnPlayerJumped}, \texttt{OnPlayerAttacked} e
\texttt{OnPlayerDashed} --- permitindo que o \texttt{ProvenanceLogger}
registre o evento sem que o \texttt{PlayerController} tenha conhecimento do
sistema de proveniência.

\subsection{Combate --- PlayerCombat e IDamageable}
\label{subsec:player-combat}

O componente \texttt{PlayerCombat} implementa a lógica de detecção e
aplicação de dano do ataque do jogador. A cada execução de ataque, o
componente utiliza o método \texttt{Physics2D.OverlapBoxAll} para detectar
todos os \textit{colliders} dentro de uma área retangular posicionada à
frente do personagem. Os objetos detectados que implementam a interface
\texttt{IDamageable} recebem dano por meio da chamada ao método
\texttt{TakeDamage(float amount)}.

A adoção da interface \texttt{IDamageable} como contrato de dano é uma
decisão de projeto que promove extensibilidade: qualquer entidade do jogo que
precise receber dano --- sejam chefes, inimigos comuns ou objetos destrutíveis
--- pode implementar essa interface sem alterar o código de ataque do jogador.
No contexto do \textit{vertical slice}, tanto o \texttt{BossHealth} (chefe
FSM) quanto o chefe PPO implementam \texttt{IDamageable}, garantindo que o
mesmo sistema de combate funcione em ambas as condições experimentais.

\subsection{Vida do Jogador --- PlayerHealth}
\label{subsec:player-health}

O componente \texttt{PlayerHealth} gerencia os pontos de vida do personagem
Elian. Ele implementa a interface \texttt{IDamageable} e mantém um valor de
vida corrente com máximo configurável (\texttt{maxHealth} = 100). Quando o
valor de vida atinge zero, o componente emite um evento de morte que é
consumido pelo \texttt{ArenaManager} para encerrar a sessão como derrota.

Esse componente foi adicionado durante a preparação para a coleta oficial,
após a identificação de que a ausência de vida real no jogador comprometeriam
a validade das métricas: sem possibilidade de derrota, o campo
\texttt{player\_damage\_taken} seria artificial e o resultado da sessão
tenderia invariavelmente a vitória.

\subsection{Boss FSM --- BossFsmController}
\label{subsec:boss-fsm-impl}

O \texttt{BossFsmController} implementa o chefe \textit{baseline}
como uma máquina de estados finitos com quatro estados, operando sobre um
\texttt{Rigidbody2D}. As transições são determinadas exclusivamente pela
distância horizontal até o jogador e pelo estado interno de
\textit{cooldown}.

A \autoref{tab:fsm-estados} descreve o comportamento associado a cada estado.

\begin{table}[htb]
    \centering
    \caption{Estados e comportamentos do chefe FSM.}
    \label{tab:fsm-estados}
    \begin{tabular}{l p{10cm}}
        \toprule
        \textbf{Estado}       & \textbf{Comportamento} \\
        \midrule
        \textit{Idle}         & O chefe permanece parado. Ocorre tipicamente
                                quando o jogador está dentro do alcance de
                                ataque, mas o \textit{cooldown} ainda não
                                expirou. \\
        \addlinespace
        \textit{Chase}        & O chefe move-se horizontalmente em direção ao
                                jogador com velocidade de perseguição
                                configurável. \\
        \addlinespace
        \textit{Attack}       & O chefe interrompe o movimento, emite o evento
                                \texttt{OnBossAttackPerformed}, exibe
                                \textit{feedback} visual e entra em
                                \textit{cooldown}. \\
        \addlinespace
        \textit{Retreat}      & O chefe move-se na direção oposta ao jogador
                                por um curto período, criando espaçamento
                                após o ataque. \\
        \bottomrule
    \end{tabular}
    \fonte{Elaborado pelo autor.}
\end{table}

Os parâmetros do FSM foram calibrados para produzir um oponente funcional mas
previsível, conforme apresentado na \autoref{tab:fsm-params}.

\begin{table}[htb]
    \centering
    \caption{Parâmetros do chefe FSM.}
    \label{tab:fsm-params}
    \begin{tabular}{l r l}
        \toprule
        \textbf{Parâmetro}          & \textbf{Valor} & \textbf{Unidade} \\
        \midrule
        \texttt{chaseSpeed}         & 2,0     & unidades/s         \\
        \texttt{attackRange}        & 1,25    & unidades           \\
        \texttt{attackDuration}     & 0,2     & segundos           \\
        \texttt{retreatDuration}    & 0,45    & segundos           \\
        \texttt{retreatSpeed}       & 1,5     & unidades/s         \\
        \texttt{attackCooldown}     & 0,8     & segundos           \\
        \bottomrule
    \end{tabular}
    \fonte{Elaborado pelo autor.}
\end{table}

A previsibilidade do FSM é uma propriedade intencional: ela garante que o
\textit{baseline} forneça um referencial estável e reprodutível para a
comparação com o agente treinado por PPO \cite{millington2019ai}.

\subsection{Dano do Boss --- BossAttackDamage}
\label{subsec:boss-attack-damage}

O componente \texttt{BossAttackDamage} é compartilhado pelas duas condições
experimentais (FSM e PPO) e aplica dano real ao jogador quando um ataque do
chefe é executado dentro do alcance. A cada evento de ataque, o componente
verifica a distância até o jogador; se esta for inferior a
\texttt{attackRange} (1,25 unidades) e o \textit{cooldown} de dano (0,8
segundos) tiver expirado, 10 pontos de dano são aplicados ao
\texttt{PlayerHealth}.

O \textit{cooldown} de dano é especialmente relevante na condição PPO: como
o agente treinado apresenta alta frequência de ações de ataque, a ausência
de \textit{cooldown} resultaria em dano instantâneo e inviabilizaria a
avaliação. O valor de 0,8 segundos é idêntico ao \texttt{attackCooldown}
do FSM, preservando a equivalência de condições.

% -----------------------------------------------------------------------------
\section{Núcleo de Proveniência}
\label{sec:nucleo-proveniencia}

O sistema de proveniência é composto por três componentes que, em conjunto,
registram, armazenam e exportam as cadeias causais de eventos de cada sessão
de jogo.

\subsection{ProvenanceGraph --- Estrutura em Memória}
\label{subsec:provenance-graph}

O \texttt{ProvenanceGraph} mantém o grafo de proveniência da sessão corrente
em memória. Cada nó do grafo é um evento com os seguintes atributos:

\begin{itemize}
    \item \texttt{eventId}: identificador único do evento;
    \item \texttt{timestamp}: marcação temporal em segundos desde o início da
          sessão;
    \item \texttt{actorId}: identificador do ator que originou o evento
          (jogador ou chefe);
    \item \texttt{actionType}: tipo da ação (e.g., \texttt{PlayerAttack},
          \texttt{BossAttack}, \texttt{BossDamageTaken});
    \item \texttt{position}: posição 2D do ator no momento do evento;
    \item \texttt{value}: valor numérico associado (e.g., quantidade de dano);
    \item \texttt{parentEventId}: referência ao evento causador.
\end{itemize}

O atributo \texttt{parentEventId} é o elemento que distingue este sistema de
um log sequencial convencional. Ao apontar para o evento que causou o evento
corrente, ele estabelece uma relação causal explícita, transformando a
sequência temporal de eventos em um grafo dirigido acíclico (DAG) com
semântica de causa e efeito \cite{provenance_base}.

\subsection{ProvenanceLogger --- Captura de Eventos}
\label{subsec:provenance-logger}

O \texttt{ProvenanceLogger} atua como ponto de integração entre o
\textit{gameplay} e o sistema de proveniência. Ele inscreve-se nos eventos
emitidos pelos componentes de jogo e, para cada evento recebido, cria um nó
correspondente no \texttt{ProvenanceGraph} com o \texttt{parentEventId}
apontando para o evento mais recente na cadeia.

Os eventos capturados pelo \texttt{ProvenanceLogger} incluem:

\begin{itemize}
    \item \textbf{Ações do jogador}: \texttt{PlayerJump},
          \texttt{PlayerAttack}, \texttt{PlayerDash} --- emitidos pelo
          \texttt{PlayerController};
    \item \textbf{Dano no jogador}: \texttt{PlayerDamageTaken} --- emitido
          pelo \texttt{PlayerHealth} ao receber dano;
    \item \textbf{Ações do chefe}: \texttt{BossAttack} --- emitido pelo
          \texttt{BossFsmController} (ou equivalente no PPO);
    \item \textbf{Dano no chefe}: \texttt{BossDamageTaken} --- emitido pelo
          \texttt{BossHealth};
    \item \textbf{Morte do chefe}: \texttt{BossDeath} --- emitido pelo
          \texttt{BossHealth} quando a vida atinge zero;
    \item \textbf{Controle de sessão}: \texttt{SessionStart} e
          \texttt{SessionEnd} --- emitidos pelo \texttt{ArenaManager}.
\end{itemize}

Um aspecto relevante de projeto é que o \texttt{ProvenanceLogger} não modifica
nem interfere no fluxo de \textit{gameplay}. Ele é um observador puro: escuta
eventos, registra nós no grafo e não emite ações de volta ao jogo. Essa
separação garante que a presença ou ausência do sistema de proveniência não
altere o comportamento da simulação, propriedade essencial para a validade da
comparação entre as condições FSM e PPO.

A \autoref{fig:cadeia-causal} ilustra um exemplo de cadeia causal típica
registrada pelo sistema.

 \begin{figure}[htb]
     \centering
     \includegraphics[width=0.8\textwidth]{imagens/cadeia-causal}
     \caption{Exemplo de cadeia causal registrada pelo sistema de
              proveniência em uma sessão de combate.}
     \label{fig:cadeia-causal}
     \fonte{Elaborado pelo autor.}
 \end{figure}

A cadeia causal típica de uma sessão vitoriosa segue o encadeamento:

\begin{quote}
    \texttt{SessionStart} $\rightarrow$ \texttt{BossAttack} $\rightarrow$
    \texttt{PlayerDamageTaken} $\rightarrow$ \texttt{PlayerAttack}
    $\rightarrow$ \texttt{PlayerDamageDealt} $\rightarrow$
    \texttt{BossDamageTaken} $\rightarrow$ \texttt{BossDeath} $\rightarrow$
    \texttt{SessionEnd}
\end{quote}

Nessa cadeia, a relação entre \texttt{PlayerAttack} e
\texttt{BossDamageTaken} não é meramente temporal: o atributo
\texttt{parentEventId} do nó \texttt{BossDamageTaken} aponta explicitamente
para o nó \texttt{PlayerAttack}, indicando causalidade direta.

\subsection{ProvenanceExporter --- Serialização em JSON}
\label{subsec:provenance-exporter}

O \texttt{ProvenanceExporter} é responsável por serializar o grafo de
proveniência em formato JSON ao término de cada sessão. A exportação é
disparada pelo evento \texttt{SessionEnd} emitido pelo
\texttt{ArenaManager}, e o arquivo resultante é gravado no diretório
\texttt{TreinamentoML/provenance\_data/} com nome baseado no identificador
da sessão.

Cada arquivo JSON contém a lista completa de nós do grafo, incluindo todos os
atributos descritos na \autoref{subsec:provenance-graph}, além de metadados da
sessão como resultado (\texttt{victory} ou \texttt{defeat}), duração e tipo
de chefe. Esse formato permite que o script Python de extração de sequências
opere diretamente sobre os arquivos exportados, sem necessidade de
transformações intermediárias.

% -----------------------------------------------------------------------------
\section{Integração com Aprendizado por Reforço}
\label{sec:integracao-rl}

A integração entre o sistema de proveniência e o aprendizado por reforço é
realizada por dois componentes complementares: o
\texttt{ProvenanceRewardShaper}, que converte as sequências vitoriosas em
sinais de recompensa, e o \texttt{BossAgent}, que implementa o agente PPO
com observações vetoriais e ações discretas.

\subsection{ProvenanceRewardShaper --- Recompensa por Proveniência}
\label{subsec:reward-shaper-impl}

O \texttt{ProvenanceRewardShaper} é o componente que materializa a conexão
entre proveniência e aprendizado por reforço. Seu funcionamento segue três
etapas:

\begin{enumerate}
    \item \textbf{Carregamento}: na inicialização, o componente lê o arquivo
          \texttt{winning\_sequences.json} --- gerado pelo script Python de
          extração --- e armazena as sequências vitoriosas em memória como
          listas de \textit{strings};
    \item \textbf{Registro}: a cada ação executada pelo agente, a
          \textit{string} correspondente à ação (e.g.,
          \texttt{"PlayerAttack"}) é inserida em um \textit{buffer} circular
          de tamanho fixo (igual ao tamanho das sequências, i.e., 3);
    \item \textbf{Comparação}: após cada inserção, o conteúdo do
          \textit{buffer} é comparado com todas as sequências vitoriosas
          carregadas. Quando há correspondência exata, o componente retorna
          uma recompensa configurável com valor padrão de $+2{,}0$.
\end{enumerate}

O \texttt{BossAgent} consulta o \texttt{ProvenanceRewardShaper} a cada passo
de decisão e, caso haja recompensa disponível, a adiciona ao sinal de
recompensa do episódio via \texttt{AddReward}. Essa recompensa auxiliar atua
em conjunto com a recompensa extrínseca do ambiente, sem substituí-la.

A decisão de isolar a lógica de \textit{Reward Shaping} em um componente
separado foi motivada por quatro razões:

\begin{enumerate}
    \item Permite testar o mecanismo de correspondência de sequências
          isoladamente, sem depender do ciclo de treinamento do ML-Agents;
    \item Mantém o código do \texttt{BossAgent} focado em observações e
          ações, sem misturar lógica de proveniência;
    \item Facilita a substituição ou extensão do esquema de \textit{Reward
          Shaping} em trabalhos futuros;
    \item Preserva a rastreabilidade da contribuição acadêmica: toda a lógica
          que conecta proveniência e RL está localizada em um único componente
          auditável.
\end{enumerate}

\subsection{Vocabulário Compartilhado de Ações}
\label{subsec:vocabulario-acoes}

Um aspecto de implementação que merece destaque é a correspondência de
vocabulário entre as ações do agente PPO e as ações humanas registradas pela
proveniência. Quando o \texttt{BossAgent} executa uma ação de pulo, ataque ou
\textit{dash}, ele registra no \texttt{ProvenanceRewardShaper} a mesma
\textit{string} utilizada pelo \texttt{ProvenanceLogger} para a ação
equivalente do jogador humano:

\begin{itemize}
    \item Ação de pulo do agente $\rightarrow$ \texttt{"PlayerJump"}
    \item Ação de ataque do agente $\rightarrow$ \texttt{"PlayerAttack"}
    \item Ação de \textit{dash} do agente $\rightarrow$ \texttt{"PlayerDash"}
\end{itemize}

Essa correspondência garante que a comparação entre o \textit{buffer} do
agente e as sequências vitoriosas extraídas dos grafos de proveniência seja
direta, sem necessidade de mapeamentos intermediários. O agente PPO, ao
executar a mesma combinação de ações que um jogador humano utilizou para
causar dano no chefe em uma sessão vitoriosa, recebe a recompensa auxiliar.

\subsection{BossAgent --- Agente PPO}
\label{subsec:boss-agent-impl}

O \texttt{BossAgent} herda da classe \texttt{Agent} do Unity ML-Agents
\cite{juliani2020unity} e implementa os três métodos fundamentais do ciclo de
treinamento: \texttt{CollectObservations}, \texttt{OnActionReceived} e
\texttt{OnEpisodeBegin}.

\subsubsection{Espaço de Observação}

O agente coleta 10 observações vetoriais a cada passo de decisão, conforme
detalhado na \autoref{tab:observacoes}.

\begin{table}[htb]
    \centering
    \caption{Observações vetoriais do \texttt{BossAgent}.}
    \label{tab:observacoes}
    \begin{tabular}{r l p{6.5cm}}
        \toprule
        \textbf{\#} & \textbf{Observação}     & \textbf{Descrição} \\
        \midrule
        1  & \texttt{bossPosition.x}           & Posição horizontal do chefe. \\
        2  & \texttt{bossPosition.y}           & Posição vertical do chefe. \\
        3  & \texttt{bossVelocity.x}           & Velocidade horizontal do chefe. \\
        4  & \texttt{bossVelocity.y}           & Velocidade vertical do chefe. \\
        5  & \texttt{relativePlayerPos.x}      & Distância relativa horizontal
                                                 até o jogador. \\
        6  & \texttt{relativePlayerPos.y}      & Distância relativa vertical
                                                 até o jogador. \\
        7  & \texttt{normalizedBossHealth}      & Vida normalizada do chefe
                                                 (intervalo $[0, 1]$). \\
        8  & \texttt{1.0} (\textit{placeholder}) & Vida do jogador.
                                                 \textit{Placeholder} fixo
                                                 utilizado durante o treino. \\
        9  & \texttt{IsGrounded()}              & Indicador binário: 1 se o
                                                 chefe está no chão, 0 caso
                                                 contrário. \\
        10 & \texttt{0.0} (\textit{placeholder}) & \textit{Cooldown} de
                                                 \textit{dash}.
                                                 \textit{Placeholder} fixo. \\
        \bottomrule
    \end{tabular}
    \fonte{Elaborado pelo autor.}
\end{table}

As observações 8 e 10 são \textit{placeholders} fixos que foram preservados
na arquitetura final do modelo. Embora não forneçam informação variável ao
agente, sua presença no vetor de observação era necessária para manter a
compatibilidade dimensional com o modelo ONNX exportado. Essa decisão é uma
consequência direta do congelamento do modelo antes da adição do
\texttt{PlayerHealth} real, conforme documentado nas decisões de escopo
(\autoref{sec:decisoes-escopo}).

\subsubsection{Espaço de Ação}

O agente opera com um espaço de ação discreto contendo seis ações mutuamente
exclusivas, apresentadas na \autoref{tab:acoes-ppo}.

\begin{table}[htb]
    \centering
    \caption{Ações discretas do \texttt{BossAgent}.}
    \label{tab:acoes-ppo}
    \begin{tabular}{r l p{7cm}}
        \toprule
        \textbf{Índice} & \textbf{Ação}          & \textbf{Efeito} \\
        \midrule
        0 & \textit{Idle}          & Nenhuma ação. O chefe permanece parado. \\
        1 & \textit{Move Left}     & Desloca o chefe para a esquerda. \\
        2 & \textit{Move Right}    & Desloca o chefe para a direita. \\
        3 & \textit{Jump}          & Aplica impulso vertical ao chefe. Registra
                                     \texttt{"PlayerJump"} no \textit{Reward
                                     Shaper}. \\
        4 & \textit{Attack Melee}  & Executa ataque corpo a corpo. Registra
                                     \texttt{"PlayerAttack"} no \textit{Reward
                                     Shaper}. \\
        5 & \textit{Dash}          & Executa \textit{dash} horizontal. Registra
                                     \texttt{"PlayerDash"} no \textit{Reward
                                     Shaper}. \\
        \bottomrule
    \end{tabular}
    \fonte{Elaborado pelo autor.}
\end{table}

As ações 3, 4 e 5 correspondem ao vocabulário de ações que o
\texttt{ProvenanceRewardShaper} monitora. As ações 0, 1 e 2 não participam
do mecanismo de \textit{Reward Shaping}, pois não fazem parte do vocabulário
de ações registrado nos grafos de proveniência.

\subsubsection{Integração do Modelo ONNX}

O \textit{run} final de treinamento (\texttt{boss\_v1\_500k}) produziu o
modelo \texttt{BossAgent.onnx} após 500.000 passos de treinamento. As
métricas finais registradas pelo TensorBoard foram: recompensa média
acumulada de $1766{,}4$ com desvio-padrão de $12{,}8$, obtidas em
aproximadamente 2.008 segundos de treinamento.

O modelo foi integrado ao projeto Unity no diretório
\texttt{Assets/Models/BossAgent.onnx} e associado ao componente
\texttt{BossAgent} no \textit{prefab} \texttt{Boss\_PPO} em modo
\texttt{InferenceOnly}. Nesse modo, o agente utiliza o modelo ONNX para
tomar decisões a cada passo, sem conexão com o treinador Python. Essa
configuração garante que o comportamento observado durante a avaliação reflita
exclusivamente a política aprendida durante o treinamento, sem qualquer ajuste
\textit{online} \cite{onnx2019}.

% -----------------------------------------------------------------------------
\section{Sistema de Métricas e Gerenciamento de Sessão}
\label{sec:sistema-metricas}

O sistema de coleta de métricas opera como uma camada de instrumentação
acoplada aos eventos já existentes no \textit{gameplay}, sem interferir na
simulação.

\subsection{ArenaManager --- Ciclo de Vida da Sessão}
\label{subsec:arena-manager}

O \texttt{ArenaManager} coordena o ciclo de vida de cada sessão de combate.
Suas responsabilidades incluem:

\begin{itemize}
    \item Emitir o evento \texttt{SessionStart} ao iniciar a cena;
    \item Detectar vitória (morte do chefe via evento do \texttt{BossHealth})
          ou derrota (morte do jogador via evento do \texttt{PlayerHealth});
    \item Emitir o evento \texttt{SessionEnd} com o resultado
          (\texttt{victory} ou \texttt{defeat});
    \item Acionar a exportação de métricas CSV e, quando habilitada, a
          exportação do grafo de proveniência em JSON.
\end{itemize}

\subsection{GameMetrics e GameMetricsExporter --- Coleta e Exportação CSV}
\label{subsec:game-metrics}

O \texttt{GameMetrics} inscreve-se nos eventos de \textit{gameplay} e
acumula contadores e métricas ao longo da sessão: duração, dano recebido pelo
chefe, dano recebido pelo jogador, contagem de ações do jogador (pulos,
ataques, \textit{dashes}) e contagem de ações do chefe. Para sessões contra o
chefe PPO, são registradas também as contagens individuais de cada uma das
seis ações discretas do \texttt{BossAgent}.

Ao término da sessão, o \texttt{GameMetricsExporter} serializa todas as
métricas em uma linha CSV e a anexa ao arquivo
\texttt{boss\_evaluation\_metrics.csv}. O exportador implementa três
cuidados de engenharia:

\begin{enumerate}
    \item \textbf{Cabeçalho estável}: o cabeçalho é escrito uma única vez, na
          criação do arquivo, garantindo que todas as linhas compartilhem a
          mesma estrutura de colunas;
    \item \textbf{Formato invariante}: valores numéricos são formatados com
          \texttt{CultureInfo.InvariantCulture}, assegurando que o separador
          decimal seja sempre o ponto, independentemente da configuração
          regional do sistema operacional;
    \item \textbf{Resolução de caminho}: o caminho do arquivo é resolvido
          relativamente à raiz do repositório, e o diretório de destino é
          criado automaticamente caso não exista.
\end{enumerate}

Essas precauções garantem que o CSV gerado seja diretamente consumível por
ferramentas de análise sem necessidade de pré-processamento.

% =============================================================================
% Capítulo 5 — Avaliação Quantitativa
% =============================================================================
\chapter{Avaliação Quantitativa}
\label{cap:avaliacao}

Este capítulo descreve o cenário, as condições e o protocolo adotados para a
avaliação quantitativa comparativa entre o chefe controlado por FSM e o chefe
controlado pelo modelo PPO treinado. O objetivo da avaliação é verificar se o
pipeline de \textit{Reward Shaping} por proveniência produziu um agente com
comportamento mensuravelmente distinto do \textit{baseline} determinístico,
utilizando métricas objetivas coletadas em sessões controladas. Os resultados
numéricos da avaliação são apresentados no \autoref{cap:resultados}.

% -----------------------------------------------------------------------------
\section{Ambiente de Avaliação}
\label{sec:ambiente-avaliacao}

A avaliação foi conduzida em duas cenas do Unity, cada uma configurada para
uma das condições experimentais. Ambas as cenas compartilham a mesma arena de
combate, os mesmos parâmetros de dano e vida, e o mesmo personagem jogável
(Elian), diferindo exclusivamente no componente que controla o chefe.

A \autoref{tab:cenas-oficiais} apresenta as cenas oficiais utilizadas.

\begin{table}[htb]
    \centering
    \caption{Cenas oficiais de avaliação e suas configurações.}
    \label{tab:cenas-oficiais}
    \begin{tabular}{l p{5cm} p{5.5cm}}
        \toprule
        \textbf{Condição} & \textbf{Cena Unity}
                          & \textbf{Configuração} \\
        \midrule
        FSM & \texttt{ArenaChefe.unity}
            & Chefe controlado por \texttt{BossFsmController}. Exporta CSV
              de métricas e JSON de proveniência. \\
        \addlinespace
        PPO & \texttt{BossArena\_PPO.unity}
            & Chefe controlado por \texttt{BossAgent} em modo
              \texttt{InferenceOnly}. Exporta CSV de métricas. \\
        \bottomrule
    \end{tabular}
    \fonte{Elaborado pelo autor.}
\end{table}

\subsection{Preparação das Cenas}
\label{subsec:preparacao-cenas}

Ambas as cenas foram preparadas para serem jogáveis pelo desenvolvedor em
condições idênticas. Essa preparação envolveu ajustes específicos para
garantir a equivalência entre as condições experimentais:

\begin{itemize}
    \item \textbf{Personagem controlável}: nas duas cenas, Elian é controlado
          pelo desenvolvedor via \textit{Input System}, com as mesmas
          mecânicas de movimento, pulo, ataque e \textit{dash};
    \item \textbf{Câmera}: o componente \texttt{HorizontalCameraFollow} foi
          adicionado para que a câmera acompanhe o personagem apenas no eixo
          horizontal, melhorando a jogabilidade sem alterar nenhum aspecto da
          IA, do sistema de recompensa, do modelo, das \textit{hitboxes} ou
          das métricas;
    \item \textbf{Parâmetros equalizados}: os valores de vida do jogador, dano
          do chefe, alcance de ataque e \textit{cooldown} de dano são
          idênticos nas duas cenas, eliminando vieses de configuração.
\end{itemize}

A cena PPO passou por uma adaptação adicional: durante o desenvolvimento, ela
era orientada ao treinamento e utilizava um \texttt{PlayerDummyAI} como
oponente simulado. Para a fase de avaliação, a dependência do
\texttt{PlayerDummyAI} foi removida e a cena foi reconfigurada para receber
entrada humana, garantindo que as sessões contra o chefe PPO seguissem o mesmo
fluxo operacional das sessões contra o chefe FSM.

\subsection{Exportação de Proveniência na Avaliação}
\label{subsec:proveniencia-avaliacao}

Uma distinção operacional entre as duas cenas é a exportação de grafos de
proveniência em JSON. Na cena FSM, a exportação ao término da sessão
(\texttt{exportOnSessionEnd}) permaneceu habilitada, gerando um arquivo JSON
por sessão. Na cena PPO, essa exportação foi desabilitada por decisão de
escopo: a métrica principal da avaliação comparativa é o CSV quantitativo, e
os JSONs da cena FSM já constituem evidência suficiente do funcionamento do
mecanismo de proveniência.

Essa decisão não afeta a validade da comparação quantitativa, pois ambas as
cenas exportam o mesmo CSV de métricas com estrutura de colunas idêntica.

% -----------------------------------------------------------------------------
\section{Condições de Vitória e Derrota}
\label{sec:condicoes-vitoria-derrota}

A definição de condições reais de vitória e derrota é um requisito
metodológico fundamental para a validade da avaliação. Uma comparação entre
chefes FSM e PPO na qual o jogador não pudesse ser derrotado produziria dados
enviesados: o campo \texttt{player\_damage\_taken} seria artificial, o
resultado de todas as sessões tenderia invariavelmente a vitória, e a
capacidade ofensiva do chefe --- que é precisamente a variável de interesse
--- não seria adequadamente capturada pelas métricas.

\subsection{Adição do Sistema de Vida Real}
\label{subsec:adicao-vida-real}

Para resolver essa limitação, o componente \texttt{PlayerHealth} foi integrado
ao personagem Elian antes do início da coleta oficial. Esse componente
implementa um sistema de pontos de vida com valor máximo de 100, dano
recebido cumulativo e emissão de evento de morte quando os pontos de vida
atingem zero. A morte do jogador é detectada pelo \texttt{ArenaManager}, que
encerra a sessão com resultado \texttt{defeat}.

Complementarmente, o componente \texttt{BossAttackDamage} foi implementado
para aplicar dano real ao jogador quando o chefe executa um ataque dentro do
alcance. Esse componente é compartilhado pelas duas condições experimentais
(FSM e PPO), garantindo que o sistema de dano opere de forma idêntica
independentemente do controlador do chefe.

\subsection{Parâmetros de Dano}
\label{subsec:parametros-dano}

A \autoref{tab:parametros-dano} apresenta os parâmetros de dano utilizados na
avaliação.

\begin{table}[htb]
    \centering
    \caption{Parâmetros de dano do chefe utilizados na avaliação.}
    \label{tab:parametros-dano}
    \begin{tabular}{l r l}
        \toprule
        \textbf{Parâmetro}                 & \textbf{Valor} & \textbf{Unidade} \\
        \midrule
        \texttt{PlayerHealth.maxHealth}    & 100   & pontos de vida \\
        \texttt{damageAmount}              & 10    & pontos por ataque \\
        \texttt{attackRange}               & 1,25  & unidades \\
        \texttt{damageCooldown}            & 0,8   & segundos \\
        \bottomrule
    \end{tabular}
    \fonte{Elaborado pelo autor.}
\end{table}

Todos os parâmetros foram mantidos idênticos nas duas condições para
eliminar vieses de configuração. A igualdade do \texttt{damageAmount} e do
\texttt{attackRange} garante que cada ataque do chefe, seja FSM ou PPO,
produza exatamente o mesmo efeito quando conectado.

\subsection{O Papel do Cooldown de Dano}
\label{subsec:cooldown-dano}

O parâmetro \texttt{damageCooldown} merece atenção especial. Esse
\textit{cooldown} define o intervalo mínimo entre duas aplicações consecutivas
de dano pelo chefe, independentemente da frequência com que o chefe executa
ações de ataque.

No contexto da condição FSM, o \textit{cooldown} de dano é pouco perceptível,
pois o \texttt{attackCooldown} interno do FSM (também de 0,8 segundos) já
limita naturalmente a frequência de ataques. Na condição PPO, contudo, o
\textit{cooldown} de dano é crucial: o agente treinado apresenta alta
frequência de ações de ataque, executando a ação \textit{Attack Melee} em
uma fração significativa dos passos de decisão. Sem o \textit{cooldown} de
dano, essa frequência elevada resultaria em múltiplas aplicações de dano
por segundo, eliminando o jogador quase instantaneamente e inviabilizando
sessões jogáveis.

O valor de 0,8 segundos foi escolhido para alinhar-se ao \textit{cooldown}
natural do FSM, garantindo que a taxa máxima de dano efetivo seja a mesma nas
duas condições. Dessa forma, a diferença na frequência de ações de ataque
entre FSM e PPO manifesta-se na contagem de ações registrada (\texttt{boss\_attack\_count}),
mas não na taxa de dano efetivo aplicado ao jogador.

\subsection{Integração com os Sistemas Existentes}
\label{subsec:integracao-sistemas}

A adição do sistema de vida real exigiu integração com três sistemas
existentes:

\begin{enumerate}
    \item \textbf{ArenaManager}: a morte do jogador (evento emitido pelo
          \texttt{PlayerHealth}) passou a ser detectada como condição de
          encerramento da sessão, com resultado \texttt{defeat};
    \item \textbf{GameMetrics}: o campo \texttt{player\_damage\_taken} passou
          a ser alimentado pelo dano real aplicado pelo
          \texttt{BossAttackDamage}, em vez de valores artificiais;
    \item \textbf{ProvenanceLogger}: o evento \texttt{PlayerDamageTaken} foi
          adicionado ao conjunto de eventos capturados, registrando cada
          instância de dano real no grafo de proveniência da sessão.
\end{enumerate}

Essas integrações seguem o mesmo padrão de comunicação por eventos descrito
no \autoref{cap:implementacao}: cada sistema inscreve-se nos eventos emitidos
pelo \texttt{PlayerHealth} sem modificar o componente emissor, preservando o
baixo acoplamento da arquitetura.

% -----------------------------------------------------------------------------
\section{Protocolo de Coleta}
\label{sec:protocolo-coleta}

O protocolo de coleta foi projetado para maximizar a comparabilidade entre as
condições experimentais dentro das restrições de escopo do trabalho.

\subsection{Procedimento}
\label{subsec:procedimento}

A coleta oficial consistiu em 20 sessões jogadas manualmente pelo
desenvolvedor, divididas em dois blocos sequenciais:

\begin{enumerate}
    \item \textbf{Bloco FSM}: 10 sessões consecutivas na cena
          \texttt{ArenaChefe.unity}, contra o chefe controlado por FSM;
    \item \textbf{Bloco PPO}: 10 sessões consecutivas na cena
          \texttt{BossArena\_PPO.unity}, contra o chefe controlado pelo
          modelo PPO em modo \texttt{InferenceOnly}.
\end{enumerate}

Em ambos os blocos, o desenvolvedor jogou normalmente como Elian, utilizando
as mecânicas disponíveis (movimento, pulo, ataque e \textit{dash}) sem
restrições ou instruções comportamentais específicas. Cada sessão encerrou-se
naturalmente por vitória (morte do chefe) ou derrota (morte do jogador), e
as métricas foram exportadas automaticamente ao término da sessão.

A \autoref{tab:amostra} sintetiza a composição da amostra oficial.

\begin{table}[htb]
    \centering
    \caption{Composição da amostra oficial da avaliação.}
    \label{tab:amostra}
    \begin{tabular}{l r}
        \toprule
        \textbf{Condição} & \textbf{Sessões} \\
        \midrule
        FSM                & 10 \\
        PPO                & 10 \\
        \midrule
        \textbf{Total}     & \textbf{20} \\
        \bottomrule
    \end{tabular}
    \fonte{Elaborado pelo autor.}
\end{table}

\subsection{Coleta Automatizada de Métricas}
\label{subsec:coleta-automatizada}

As métricas de cada sessão foram registradas automaticamente pelo sistema
\texttt{GameMetrics} e exportadas pelo \texttt{GameMetricsExporter} para o
arquivo CSV \texttt{boss\_evaluation\_metrics.csv}. Esse processo não requer
intervenção manual do jogador: os contadores são inicializados
automaticamente ao início da sessão e a linha CSV é gravada ao término,
garantindo a lisura dos dados coletados.

O cabeçalho do CSV contém 18 campos, abrangendo identificação da sessão
(\texttt{session\_id}, \texttt{boss\_type}), resultado
(\texttt{result}), temporização (\texttt{start\_time\_seconds},
\texttt{end\_time\_seconds}, \texttt{duration\_seconds}), métricas de
combate (\texttt{boss\_damage\_taken}, \texttt{player\_damage\_taken}) e
contagens detalhadas de ações do jogador e do chefe. Para sessões contra o
chefe PPO, seis campos adicionais registram a distribuição das ações
discretas do agente (\texttt{ppo\_idle\_count} a \texttt{ppo\_dash\_count}).

Três cuidados de engenharia garantem a confiabilidade do CSV exportado:

\begin{enumerate}
    \item O cabeçalho é escrito uma única vez, na criação do arquivo,
          assegurando estrutura de colunas consistente;
    \item Valores numéricos são formatados com
          \texttt{CultureInfo.InvariantCulture}, de modo que o separador
          decimal seja sempre o ponto, independentemente da configuração
          regional do sistema;
    \item O caminho do arquivo é resolvido relativamente à raiz do
          repositório, e o diretório de destino é criado automaticamente
          caso não exista.
\end{enumerate}

\subsection{Controle de Linhas Extras}
\label{subsec:controle-linhas}

Após a coleta oficial das 20 sessões, sessões adicionais foram realizadas
para fins de gravação de evidências visuais (vídeos demonstrativos do
comportamento dos chefes FSM e PPO). Essas sessões adicionais geraram linhas
extras no CSV, que não fazem parte da amostra oficial.

Para evitar contaminação dos dados, a análise apresentada no
\autoref{cap:resultados} considera exclusivamente as primeiras 20 linhas do
CSV, correspondentes às 10 sessões FSM e 10 sessões PPO da coleta oficial.
Qualquer linha posterior é descartada da análise.

\subsection{Considerações sobre Validade}
\label{subsec:validade}

O protocolo adotado apresenta características que favorecem a validade interna
da comparação:

\begin{itemize}
    \item \textbf{Condições equalizadas}: parâmetros de dano, vida e alcance
          são idênticos nas duas cenas, garantindo que diferenças nas
          métricas reflitam exclusivamente o comportamento do chefe;
    \item \textbf{Coleta automatizada}: a exportação de métricas não depende
          de registro manual, eliminando erros de transcrição;
    \item \textbf{Modelo congelado}: o chefe PPO opera em inferência pura, sem
          ajustes \textit{online}, garantindo reprodutibilidade do
          comportamento;
    \item \textbf{Mesmo jogador}: todas as sessões foram conduzidas pelo
          mesmo operador, eliminando variância interpessoal.
\end{itemize}

As limitações do protocolo --- em particular, a amostra reduzida, a avaliação
conduzida pelo próprio desenvolvedor e a ausência de significância estatística
formal --- são reconhecidas e discutidas no \autoref{cap:discussao}.

% =============================================================================
% Capítulo 6 — Resultados
% =============================================================================
\chapter{Resultados}
\label{cap:resultados}

Este capítulo apresenta os resultados quantitativos obtidos na avaliação
comparativa entre os chefes FSM e PPO. Os dados são expostos de forma
descritiva, sem interpretação ou análise de causalidade --- a discussão dos
resultados é realizada no \autoref{cap:discussao}.

% -----------------------------------------------------------------------------
\section{Convergência do Treinamento}
\label{sec:convergencia}

O treinamento do agente PPO foi executado no \textit{run}
\texttt{boss\_v1\_500k}, que atingiu o limite configurado de 500.000 passos de
treinamento em aproximadamente 2.008 segundos (cerca de 33 minutos). A
\autoref{tab:metricas-treino} apresenta as métricas finais registradas pelo
TensorBoard.

\begin{table}[htb]
    \centering
    \caption{Métricas finais do treinamento PPO (\textit{run}
             \texttt{boss\_v1\_500k}).}
    \label{tab:metricas-treino}
    \begin{tabular}{l r}
        \toprule
        \textbf{Métrica}                 & \textbf{Valor} \\
        \midrule
        Recompensa média acumulada       & $1\,766{,}4$ \\
        Desvio-padrão da recompensa      & $12{,}8$ \\
        Tempo total de treinamento       & $2\,007{,}5$ s \\
        Passos totais                    & $500\,000$ \\
        \bottomrule
    \end{tabular}
    \fonte{Dados coletados via TensorBoard.}
\end{table}

O baixo desvio-padrão da recompensa ($12{,}8$) em relação à média
($1\,766{,}4$) indica que, nos episódios finais do treinamento, o agente
apresentava desempenho consistente, conforme ilustrado na evolução da recompensa acumulada (\autoref{fig:cumulative-reward}). Além disso, a convergência do modelo também é evidenciada pela estabilização do erro de valor (\textit{Value Loss}), demonstrando que a rede neural aprendeu a prever as recompensas do ambiente com precisão ao longo dos 500 mil passos (\autoref{fig:tb-loss}).

 \begin{figure}[htb]
     \centering
     \includegraphics[width=0.85\textwidth]{imagens/Environment-Cumulative Reward}
     \caption{Evolução da recompensa acumulada ao longo do treinamento
              (\textit{Environment/Cumulative Reward}).}
     \label{fig:cumulative-reward}
     \fonte{Captura do TensorBoard.}
 \end{figure}

 \begin{figure}[htb]
    \centering
    \includegraphics[width=0.85\textwidth]{imagens/Losses-Value Loss.png}
    \caption{Evolução do erro de valor (\textit{Value Loss}) indicando a estabilização do aprendizado da rede neural.}
    \label{fig:tb-loss}
    \fonte{Captura do TensorBoard.}
\end{figure}

% TODO: Inserir gráfico de Episode Length
% \begin{figure}[htb]
%     \centering
%     \includegraphics[width=0.85\textwidth]{Environment-Episode Length}
%     \caption{Evolução do comprimento dos episódios ao longo do treinamento
%              (\textit{Environment/Episode Length}).}
%     \label{fig:episode-length}
%     \fonte{Captura do TensorBoard.}
% \end{figure}

% -----------------------------------------------------------------------------
\section{Comparativo Quantitativo: FSM vs.\ PPO}
\label{sec:comparativo}

Esta seção apresenta os dados numéricos das 20 sessões oficiais da
avaliação: 10 contra o chefe FSM e 10 contra o chefe PPO.

\subsection{Resultados Gerais}
\label{subsec:resultados-gerais}

A \autoref{tab:resultados-gerais} sintetiza os resultados globais por
condição experimental.

\begin{table}[htb]
    \centering
    \caption{Resultados gerais da avaliação por condição experimental.}
    \label{tab:resultados-gerais}
    \begin{tabular}{l r r r r}
        \toprule
        \textbf{Condição}
            & \textbf{Sessões}
            & \textbf{Vitórias}
            & \textbf{Derrotas}
            & \textbf{Taxa de vitória} \\
        \midrule
        FSM & 10 & 10 & 0 & 100\% \\
        PPO & 10 &  9 & 1 &  90\% \\
        \bottomrule
    \end{tabular}
    \fonte{Dados da coleta oficial.}
\end{table}

Na condição FSM, o jogador venceu todas as 10 sessões. Na condição PPO, o
jogador venceu 9 das 10 sessões, com uma derrota registrada --- a única
ocorrência em toda a amostra de 20 sessões em que o chefe eliminou o jogador.

\subsection{Métricas de Combate}
\label{subsec:metricas-combate}

A \autoref{tab:metricas-comparativas} apresenta as médias e desvios-padrão
das métricas de combate para cada condição.

\begin{table}[htb]
    \centering
    \caption{Médias e desvios-padrão das métricas de combate por condição
             (média $\pm$ desvio-padrão, $n = 10$ por condição).}
    \label{tab:metricas-comparativas}
    \begin{tabular}{l r r}
        \toprule
        \textbf{Métrica}
            & \textbf{FSM}
            & \textbf{PPO} \\
        \midrule
        Duração da luta (s)
            & $5{,}97 \pm 2{,}79$
            & $15{,}18 \pm 7{,}74$ \\
        \addlinespace
        Dano recebido pelo chefe
            & $100{,}00 \pm 0{,}00$
            & $95{,}00 \pm 15{,}81$ \\
        \addlinespace
        Dano recebido pelo jogador
            & $50{,}00 \pm 14{,}14$
            & $56{,}00 \pm 22{,}21$ \\
        \addlinespace
        Pulos do jogador
            & $0{,}30 \pm 0{,}48$
            & $3{,}50 \pm 3{,}44$ \\
        \addlinespace
        Ataques do jogador
            & $10{,}70 \pm 0{,}95$
            & $15{,}80 \pm 2{,}57$ \\
        \addlinespace
        \textit{Dashes} do jogador
            & $0{,}20 \pm 0{,}63$
            & $0{,}50 \pm 1{,}08$ \\
        \addlinespace
        Ataques do chefe
            & $5{,}60 \pm 1{,}58$
            & $296{,}60 \pm 186{,}86$ \\
        \bottomrule
    \end{tabular}
    \fonte{Dados da coleta oficial (20 sessões).}
\end{table}

Observa-se que a duração média das sessões na condição PPO ($15{,}18$ s) é
aproximadamente 2,5 vezes superior à da condição FSM ($5{,}97$ s). O dano
recebido pelo chefe foi de $100{,}00$ em todas as sessões FSM, enquanto na
condição PPO a média foi de $95{,}00$, refletindo a sessão em que o jogador
foi derrotado antes de eliminar o chefe. O dano recebido pelo jogador foi
ligeiramente superior na condição PPO ($56{,}00$) em relação à FSM
($50{,}00$).

A contagem de ataques do jogador foi maior na condição PPO ($15{,}80$) do que
na FSM ($10{,}70$), assim como a contagem de pulos ($3{,}50$ vs.\ $0{,}30$) e
\textit{dashes} ($0{,}50$ vs.\ $0{,}20$).

A diferença mais acentuada entre as duas condições reside na contagem de
ataques do chefe: $5{,}60 \pm 1{,}58$ na condição FSM contra $296{,}60 \pm
186{,}86$ na condição PPO.

\subsection{Distribuição de Ações do Agente PPO}
\label{subsec:distribuicao-acoes}

Para a condição PPO, o sistema de métricas registrou a contagem individual de
cada uma das seis ações discretas do \texttt{BossAgent} ao longo das 10
sessões oficiais. A \autoref{tab:acoes-ppo} apresenta esses dados.

\begin{table}[htb]
    \centering
    \caption{Distribuição das ações discretas do agente PPO nas 10 sessões
             oficiais.}
    \label{tab:acoes-ppo}
    \begin{tabular}{l r r r}
        \toprule
        \textbf{Ação}
            & \textbf{Total}
            & \textbf{Média $\pm$ DP}
            & \textbf{\% aprox.} \\
        \midrule
        \textit{Attack}
            & 2\,966  & $296{,}60 \pm 186{,}86$ & 39,05\% \\
        \textit{Jump}
            & 2\,123  & $212{,}30 \pm 149{,}10$ & 27,95\% \\
        \textit{Dash}
            &    796  & $79{,}60 \pm 40{,}77$   & 10,48\% \\
        \textit{Move Right}
            &    698  & $69{,}80 \pm 31{,}94$   &  9,19\% \\
        \textit{Move Left}
            &    550  & $55{,}00 \pm 31{,}67$   &  7,24\% \\
        \textit{Idle}
            &    463  & $46{,}30 \pm 25{,}38$   &  6,10\% \\
        \midrule
        \textbf{Total}
            & \textbf{7\,596} & --- & 100\% \\
        \bottomrule
    \end{tabular}
    \fonte{Dados da coleta oficial.}
\end{table}

As ações \textit{Attack} e \textit{Jump} correspondem, juntas, a
aproximadamente 67\% de todas as ações executadas pelo agente. As ações de
movimentação (\textit{Move Left} e \textit{Move Right}) somam 16,43\%, e o
\textit{Idle} representa 6,10\% das decisões. As três ações que participam
do mecanismo de \textit{Reward Shaping} --- \textit{Attack}, \textit{Jump} e
\textit{Dash} --- totalizam 77,48\% das ações executadas.

% =============================================================================
% Capítulo 7 — Discussão
% =============================================================================
\chapter{Discussão}
\label{cap:discussao}

Este capítulo interpreta os resultados apresentados no
\autoref{cap:resultados}, analisando o comportamento exibido pelo agente PPO
em comparação com o \textit{baseline} FSM, discutindo as limitações
observadas e avaliando em que medida a proposta central do trabalho ---
utilizar proveniência como fonte de \textit{Reward Shaping} para aprendizado
por reforço --- foi validada pelos dados coletados.

% -----------------------------------------------------------------------------
\section{Análise de Comportamento e Agressividade}
\label{sec:analise-comportamento}

Os dados quantitativos da avaliação revelam um perfil comportamental
marcadamente distinto entre as duas condições experimentais. A análise a
seguir examina as dimensões mais relevantes dessa distinção.

\subsection{Duração e Persistência}
\label{subsec:duracao-persistencia}

A duração média das sessões contra o chefe PPO ($15{,}18 \pm 7{,}74$ s) é
aproximadamente 2,5 vezes superior à duração média contra o chefe FSM
($5{,}97 \pm 2{,}79$ s). Esse resultado indica que o chefe PPO é capaz de
prolongar o combate de forma consistente, resistindo por mais tempo antes de
ser derrotado. A maior variabilidade observada na condição PPO (desvio-padrão
de $7{,}74$ s, contra $2{,}79$ s no FSM) reflete a natureza menos previsível
do agente treinado: enquanto o FSM segue padrões fixos que conduzem a durações
relativamente uniformes, o PPO produz sessões com dinâmicas variadas, oscilando
entre 8,8 e 36,1 segundos na amostra oficial.

A persistência do chefe PPO não decorre de uma estratégia defensiva
elaborada --- os dados mostram que o agente é predominantemente ofensivo ---,
mas sim da pressão constante que exerce sobre o jogador. Ao atacar e se
movimentar com alta frequência, o agente dificulta a aproximação do jogador
e impõe um ritmo de combate mais intenso, o que prolonga naturalmente a
duração das sessões.

\subsection{Agressividade Ofensiva}
\label{subsec:agressividade}

A diferença mais expressiva entre as duas condições reside na contagem de
ataques do chefe. O chefe FSM executou, em média, $5{,}60 \pm 1{,}58$ ataques
por sessão, enquanto o chefe PPO executou $296{,}60 \pm 186{,}86$ --- uma
diferença de aproximadamente \textbf{53 vezes}. A distribuição de ações do
agente PPO confirma essa postura ofensiva: 39,05\% de todas as decisões do
agente correspondem à ação \textit{Attack}, seguida por \textit{Jump} com
27,95\%. Juntas, essas duas ações representam aproximadamente 67\% do
comportamento total do agente.

Essa agressividade reflete-se no comportamento adaptativo do jogador. Nas
sessões contra o PPO, o jogador executou mais pulos ($3{,}50$ vs.\ $0{,}30$),
mais ataques ($15{,}80$ vs.\ $10{,}70$) e mais \textit{dashes} ($0{,}50$ vs.\
$0{,}20$), sugerindo que a pressão exercida pelo chefe PPO forçou uma postura
mais ativa e evasiva por parte do jogador. A necessidade de mais ataques para
vencer indica que o jogador precisou investir mais esforço ofensivo para
superar o chefe PPO, ao passo que contra o FSM bastavam estratégias mais
diretas e econômicas.

\subsection{Capacidade de Ameaça}
\label{subsec:capacidade-ameaca}

O dano médio recebido pelo jogador foi ligeiramente superior na condição PPO
($56{,}00 \pm 22{,}21$) em comparação com a condição FSM ($50{,}00 \pm
14{,}14$). Embora a diferença nas médias seja moderada --- em parte limitada
pelo \textit{cooldown} de dano que equaliza a taxa de dano efetivo entre as
condições ---, o dado qualitativamente mais significativo é a ocorrência da
\textbf{única derrota do jogador} em toda a amostra de 20 sessões.

Nessa sessão específica, o jogador recebeu 100 pontos de dano (totalidade de
sua vida) e o chefe sobreviveu com 50 pontos de vida remanescentes. Essa
derrota demonstra que o agente PPO é capaz de produzir cenários de combate nos
quais o chefe representa uma ameaça real à sobrevivência do jogador, ao
contrário do FSM, cuja previsibilidade permitiu ao jogador vencer todas as
sessões sem exceção.

% -----------------------------------------------------------------------------
\section{O Fenômeno da Repetição Mecânica}
\label{sec:repeticao-mecanica}

A despeito dos indicadores positivos de agressividade e persistência, o
comportamento do agente PPO apresenta uma limitação que deve ser discutida com
transparência: a alta repetição de ações, comumente referida como
\textit{spam}.

\subsection{Caracterização do Fenômeno}
\label{subsec:caracterizacao-spam}

Os dados de distribuição de ações revelam que o agente concentra 39,05\%
de suas decisões na ação \textit{Attack} e 27,95\% na ação \textit{Jump}.
Essa concentração produz um comportamento visualmente repetitivo: o chefe PPO
ataca e pula de forma quase contínua, com pouca variação tática aparente. A
consequência visual é que o chefe pode parecer ``preso'' na animação de
ataque, pois o \textit{feedback} visual representa a última ação executada e
o agente ataca com altíssima frequência.

Embora o \textit{cooldown} de dano (0,8 segundos) implementado no
\texttt{BossAttackDamage} limite a taxa de dano efetivo --- impedindo que o
\textit{spam} resulte em morte instantânea do jogador ---, a contagem de ações
registrada no CSV reflete todas as \textit{tentativas} de ataque, não apenas
aquelas que efetivamente aplicaram dano. Assim, a grande maioria dos 296,60
ataques médios por sessão não gera dano real; eles representam decisões do
agente que são neutralizadas pelo mecanismo de \textit{cooldown}.

\subsection{Origem do Comportamento}
\label{subsec:origem-comportamento}

A repetição mecânica pode ser atribuída à combinação de dois fatores
estruturais:

\begin{enumerate}
    \item \textbf{Ausência de custo no espaço de ações}: o espaço de ações
          discretas do \texttt{BossAgent} não incorpora \textit{cooldowns}
          internos nem penalidades por repetição. Cada uma das seis ações pode
          ser executada em qualquer passo de decisão, sem restrição temporal
          ou custo energético. O algoritmo PPO, ao identificar que ações de
          ataque geram recompensa --- tanto extrínseca (por causar dano)
          quanto auxiliar (via \textit{Reward Shaping}) ---, convergiu para
          uma política que maximiza a frequência dessas ações;

    \item \textbf{Modelo congelado sem pós-processamento}: o modelo avaliado
          é o resultado direto do \textit{run} de treinamento, exportado como
          ONNX e integrado sem ajustes. A decisão de preservar o modelo
          congelado é metodologicamente correta --- a avaliação deve medir o
          comportamento aprendido, não um comportamento corrigido
          \textit{post hoc} --- mas implica que qualquer tendência degenerada
          da política aprendida é transportada integralmente para a avaliação.
\end{enumerate}

Essa observação sugere que, em trabalhos futuros, a inclusão de
\textit{cooldowns} diretamente no espaço de ações do agente ou a adição de
penalidades por repetição na função de recompensa poderiam mitigar o fenômeno
sem comprometer a autonomia do PPO.

\subsection{Preservação Metodológica}
\label{subsec:preservacao-metodologica}

É importante registrar que a decisão de não corrigir o comportamento
repetitivo antes da análise foi deliberada. Corrigir o \textit{spam} ---
por exemplo, adicionando um \textit{cooldown} forçado às ações do agente
antes da coleta --- descaracterizaria a avaliação, pois introduziria uma
modificação \textit{post hoc} que alteraria os dados coletados sem refletir
uma melhoria genuína no aprendizado. O comportamento repetitivo é um
resultado experimental do modelo congelado e, como tal, deve ser reportado
e discutido, não ocultado.

% -----------------------------------------------------------------------------
\section{Validação da Proposta: Reward Shaping por Proveniência}
\label{sec:validacao-proposta}

A questão central do trabalho é se sequências extraídas de grafos de
proveniência podem ser utilizadas como sinal auxiliar de recompensa para
orientar o treinamento de um agente PPO, produzindo um comportamento
mensuravelmente distinto do \textit{baseline} FSM --- sem recorrer a
imitação direta ou demonstrações gravadas.

\subsection{Evidências de Distinção Comportamental}
\label{subsec:evidencias-distincao}

Os dados da avaliação fornecem evidências consistentes de que o objetivo foi
alcançado. O agente PPO difere do \textit{baseline} FSM em múltiplas
dimensões mensuráveis:

\begin{itemize}
    \item \textbf{Duração}: sessões 2,5 vezes mais longas;
    \item \textbf{Volume de ações}: 53 vezes mais ataques por sessão;
    \item \textbf{Ameaça}: única condição a produzir derrota do jogador;
    \item \textbf{Distribuição de ações}: política concentrada em ações
          ofensivas, contrastando com o padrão cíclico regular do FSM;
    \item \textbf{Impacto no jogador}: aumento de pulos, ataques e
          \textit{dashes} do jogador na condição PPO, indicando adaptação a
          um oponente mais exigente.
\end{itemize}

Essas diferenças não são marginais nem atribuíveis a ruído amostral: a
magnitude das diferenças --- particularmente na contagem de ataques e na
duração --- excede largamente os desvios-padrão de ambas as condições.

\subsection{O Papel da Proveniência}
\label{subsec:papel-proveniencia}

O mecanismo de \textit{Reward Shaping} por proveniência operou conforme
projetado. As três ações que participam do \textit{Reward Shaping} ---
\textit{Attack}, \textit{Jump} e \textit{Dash} --- totalizam 77,48\% de
todas as ações executadas pelo agente, indicando que o sinal auxiliar de
recompensa exerceu influência mensurável sobre a distribuição da política
aprendida. O agente desenvolveu preferência pelas ações que eram
recompensadas quando executadas em sequências compatíveis com padrões
vitoriosos extraídos dos grafos de proveniência.

Esse resultado é significativo porque demonstra que a proveniência cumpriu
sua função de \textbf{fonte de informação causal} para o treinamento: as
sequências extraídas dos grafos --- que codificam padrões de ações associados
a desfechos favoráveis --- foram efetivamente incorporadas ao comportamento
do agente por meio do \textit{Reward Shaping}, sem a necessidade de
demonstrações diretas, arquivos \texttt{.demo} ou discriminadores adversários
(GAIL).

\subsection{Distinção em Relação à Imitação Direta}
\label{subsec:distincao-imitacao}

É fundamental distinguir o resultado obtido de um resultado de imitação. O
agente PPO \textbf{não copiou} o comportamento do jogador humano. Ele
explorou o ambiente livremente via PPO e recebeu recompensa adicional quando
executou sequências de ações compatíveis com padrões vitoriosos. O
comportamento resultante --- altamente agressivo e repetitivo --- não se
assemelha ao comportamento do jogador humano nas sessões originais de coleta
de proveniência. Isso confirma que o agente desenvolveu sua própria política
a partir dos sinais de recompensa, não uma réplica do comportamento observado.

Essa distinção é relevante tanto do ponto de vista técnico --- o pipeline não
utiliza GAIL \cite{ho2016gail} nem demonstrações gravadas --- quanto do ponto
de vista da contribuição acadêmica: o trabalho demonstra uma via alternativa
para orientar o aprendizado de agentes de jogo, baseada em informação causal
rastreável.

\subsection{Qualificações Necessárias}
\label{subsec:qualificacoes}

Os resultados devem ser interpretados com as seguintes qualificações, que
delimitam o alcance da contribuição:

\begin{enumerate}
    \item A amostra de 20 sessões (10 por condição) permite identificar
          diferenças comportamentais observadas, mas não estabelece
          significância estatística formal;
    \item A avaliação foi conduzida pelo próprio desenvolvedor, que conhece
          os padrões de ambos os chefes, introduzindo potencial viés de
          familiaridade;
    \item O modelo PPO não foi retreinado após a adição do
          \texttt{PlayerHealth} real; as observações de vida do jogador e
          \textit{cooldown} de \textit{dash} permaneceram como
          \textit{placeholders} fixos;
    \item A cena PPO não exportou grafos de proveniência durante a avaliação,
          por decisão de escopo;
    \item Os dados documentam um comportamento \textit{distinto}, não
          necessariamente \textit{superior}: as métricas provam diferença
          comportamental observada, não superioridade universal do agente PPO
          sobre o FSM.
\end{enumerate}

\subsection{Síntese}
\label{subsec:sintese-discussao}

Em síntese, o trabalho demonstrou que sequências extraídas de grafos de
proveniência podem ser utilizadas como sinal auxiliar de recompensa em um
agente PPO, produzindo um chefe com comportamento mensuravelmente distinto
do \textit{baseline} FSM em um ambiente de jogo 2D. O agente resultante é
mais agressivo, mais persistente e capaz de ameaçar a sobrevivência do
jogador, embora o comportamento aprendido apresente repetição mecânica que
compromete a percepção de sofisticação tática. A contribuição reside na
demonstração de viabilidade da integração entre proveniência e aprendizado
por reforço em um protótipo jogável, abrindo caminhos para refinamentos
futuros.

% -----------------------------------------------------------------------------
\section{Trabalhos Futuros}
\label{sec:trabalhos-futuros}

As limitações identificadas indicam direções concretas para trabalhos futuros:

\begin{enumerate}
    \item \textbf{Cooldowns no espaço de ações}: incorporar
          \textit{cooldowns} diretamente às ações discretas do agente,
          impedindo a repetição imediata e incentivando maior variedade
          tática;
    \item \textbf{Retreino com observações completas}: retreinar o modelo
          com a vida real do jogador e o \textit{cooldown} de \textit{dash}
          como observações efetivas, permitindo que o agente adapte seu
          comportamento ao estado do oponente;
    \item \textbf{\textit{Curriculum learning}}: introduzir fases de
          treinamento com dificuldade crescente, permitindo que o agente
          desenvolva estratégias progressivamente mais sofisticadas;
    \item \textbf{Avaliação com playtesters externos}: conduzir avaliação
          com jogadores externos e instrumentos de percepção subjetiva
          (questionários), complementando as métricas quantitativas com
          dados de experiência do jogador;
    \item \textbf{\textit{Self-play}}: investigar o uso de \textit{self-play}
          como mecanismo complementar de treinamento, permitindo que o agente
          evolua contra versões anteriores de si mesmo;
    \item \textbf{\textit{Reward Shaping} adaptativo}: explorar esquemas de
          \textit{Reward Shaping} que atualizem as sequências-alvo ao longo
          do treinamento, à medida que novas sessões de proveniência são
          coletadas.
\end{enumerate}

% =============================================================================
% Capítulo 8 — Conclusão
% =============================================================================
\chapter{Conclusão}
\label{cap:conclusao}

A criação de comportamentos de NPCs em jogos digitais enfrenta uma tensão
recorrente entre previsibilidade e adaptabilidade. Máquinas de Estados Finitos
oferecem controle total e facilidade de depuração, mas produzem oponentes cujos
padrões se tornam evidentes após poucas interações. O Aprendizado por Reforço
permite a emergência de comportamentos adaptativos, mas impõe o desafio de
definir funções de recompensa que orientem o agente sem recorrer a modelagem
manual exaustiva ou à imitação direta de demonstrações humanas, que vincula o
agente ao comportamento observado sem preservar a rastreabilidade causal das
decisões de \textit{design}.

O presente trabalho endereçou essa tensão por meio de uma abordagem que
integra Proveniência de Dados e Aprendizado por Reforço. Foi desenvolvido o
\textit{vertical slice} do jogo \textit{Chronicles Of The Lost Word}, um
protótipo 2D em Unity centrado em uma arena de combate contra um chefe. O
pipeline implementado percorreu todas as etapas propostas: sessões jogadas
contra o chefe FSM geraram grafos de proveniência com cadeias causais de
eventos; um processo automatizado extraiu sequências de ações associadas a
desfechos vitoriosos; essas sequências foram carregadas por um módulo de
\textit{Reward Shaping} que concedeu recompensa adicional ao agente PPO
durante o treinamento via Unity ML-Agents; e o modelo resultante foi exportado
em formato ONNX e integrado ao jogo em modo de inferência para avaliação
quantitativa.

Os resultados da avaliação --- 10 sessões contra o chefe FSM e 10 contra o
chefe PPO, com métricas coletadas automaticamente --- confirmaram que a
metodologia produziu um agente com comportamento mensuravelmente distinto do
\textit{baseline}. O chefe PPO gerou sessões de combate 2,5 vezes mais longas,
executou 53 vezes mais ações de ataque por sessão e foi a única condição
experimental a produzir uma derrota do jogador na amostra. As ações
participantes do mecanismo de \textit{Reward Shaping} --- ataque, pulo e
\textit{dash} --- corresponderam a 77,48\% das decisões do agente, indicando
que o sinal auxiliar derivado da proveniência exerceu influência direta sobre
a política aprendida.

Esses resultados validam a tese central do trabalho: sequências extraídas de
grafos de proveniência podem ser utilizadas como sinal auxiliar de recompensa
em um agente PPO, orientando o aprendizado sem imitação direta e produzindo
um comportamento que, embora distinto do humano, é rastreável às decisões
causais que o originaram. O agente não reproduziu trajetórias gravadas nem
utilizou discriminadores adversários; ele explorou o ambiente livremente e
convergiu para uma política influenciada por padrões causais vitoriosos ---
uma via alternativa ao GAIL \cite{ho2016gail} que preserva a interpretabilidade
do processo de treinamento.

As limitações observadas --- em particular, a repetição mecânica de ações
decorrente da ausência de \textit{cooldowns} no espaço de ações do agente e
do modelo congelado sem retreino após ajustes na avaliação --- não invalidam
a contribuição; elas a contextualizam. O comportamento repetitivo é uma
consequência documentada e explicável de escolhas de projeto específicas, e
sua correção constitui um passo natural de refinamento, não uma refutação da
abordagem. A amostra reduzida e a avaliação conduzida pelo próprio
desenvolvedor delimitam o alcance das conclusões à demonstração de viabilidade,
sem pretensão de generalização estatística.

Em última análise, o trabalho demonstrou que a proveniência de dados ---
originalmente concebida para rastreamento de linhagem e análise
\textit{post-mortem} de sessões --- pode ser reaproveitada como ponte entre
o \textit{gameplay} humano e o treinamento de agentes inteligentes em jogos.
Essa integração abre um espaço de investigação no qual grafos causais de
sessões de jogo deixam de ser apenas ferramentas de análise e passam a atuar
como fontes ativas de orientação para o aprendizado, combinando a
rastreabilidade da proveniência com a capacidade adaptativa do aprendizado
por reforço.


% ===========================
% ELEMENTOS PÓS-TEXTUAIS
% ===========================
\postextual

% --- Referências bibliográficas ---
\bibliography{referencias}

% --- Apêndices (opcional) ---
% \begin{apendicesenv}
%     \partapendices
%     \input{apendices/apendice-a}
% \end{apendicesenv}

% --- Anexos (opcional) ---
% \begin{anexosenv}
%     \partanexos
%     \input{anexos/anexo-a}
% \end{anexosenv}

\end{document}
